%% 8장 볼록성과 최적성 (새 목차, 2026-11-25)
\chapter{Convexity and Optimality}
\label{ch:convexity}

\abstract{최대우도 추정, 최소제곱, 벌점 회귀는 모두 최솟값을 찾는 문제이고 계산이 끝났다는 판단은 최적성 조건의 확인이다. 이 장은 극소점의 조건, 볼록 집합과 볼록 함수, 부분기울기, KKT 조건과 Lagrange 쌍대를 세운다. 구현으로는 stardust의 사영 기울기 판정, 비음 최소제곱의 활성 집합법, 라소의 쌍대 간격을 따라가고 유한 정밀도가 최솟값의 위치에 주는 한계를 잰다.}

함수의 최솟값을 찾는 알고리즘은 결국 언제 멈출지를 정해야 한다. 기울기가 $0$에 가까우면 멈추는 규칙은 자연스러워 보인다. 그러나 그 점이 정말 최소점인지, 제약이 있으면 무엇이 기울기의 자리를 대신하는지, $10^{-8}$이라는 기울기가 함수값으로 얼마나 최적에 가깝다는 뜻인지는 따로 따져야 한다. 이 질문들의 답이 최적성 조건과 쌍대성의 이론이다. 볼록성이 있으면 이 이론은 특히 강해진다. 국소 최소와 전역 최소가 같아진다. 또 쌍대 문제의 값이 원래 문제의 값을 아래에서 막아 주므로 계산한 답이 최적에서 얼마나 떨어져 있는지를 계산 가능한 수 하나로 증명할 수 있다. 이 장은 그 이론을 세우고 stardust의 멈춤 판정과 활성 집합 알고리즘, 쌍대 간격 계산이 이론을 어떻게 옮겼는지 따라간다. 선 탐색, 준뉴턴법, 신뢰 영역 같은 최적화 알고리즘 자체는 9권에서 다룬다.

\section{Minima and Optimality Conditions}
\label{sec:optimality}

\subsection{Existence of Minimizers}

함수 $f:\mathbb R^n\to\mathbb R$에 대해 $x_\ast$의 어떤 근방에서 $f(x_\ast)\le f(x)$이면 $x_\ast$를 \term{국소 최소점}{local minimizer}, 모든 $x$에서 그러면 \term{전역 최소점}{global minimizer}이라 한다. 근방의 $x\ne x_\ast$에서 부등식이 엄격하면 \term{엄격한}{strict} 최소점이다. 최소점이 늘 있지는 않다. $e^x$는 아래로 유계이지만 최솟값이 없다. $x\sin(1/x)$처럼 최소점이 끝없이 쌓이는 함수도 있다.

\begin{proposition}
\label{prop:existence}
$f$가 연속이고 어떤 $x_0$에서 준위 집합 $\{x:f(x)\le f(x_0)\}$이 유계이면 $f$는 전역 최소점을 갖는다. 특히 $\|x\|\to\infty$일 때 $f(x)\to\infty$인 \term{강제적}{coercive} 함수는 전역 최소점을 갖는다.
\end{proposition}

\begin{proof}
준위 집합은 연속 함수의 닫힌 집합의 역상이므로 닫혀 있고 유계이므로 컴팩트하다. Weierstrass 정리에서 $f$는 그 위에서 최솟값을 갖고 준위 집합 밖에서는 $f>f(x_0)$이므로 그 점이 전역 최소점이다. 강제적이면 $\|x\|>\rho$에서 $f(x)>f(x_0)$인 $\rho$가 있어 준위 집합이 공 안에 들어간다.
\end{proof}

최소제곱 $\|Ax-b\|^2$은 $A$의 열이 독립이면 강제적이고 그렇지 않으면 영공간 방향으로 평평하지만 여전히 최솟값을 갖는다. 로지스틱 회귀의 음의 로그 우도는 자료가 완전히 분리되면 아래로 유계이면서 최소점이 없어 계수가 무한대로 달아난다. 이런 경우를 알아채고 알리는 일도 구현이 맡는다(12권).

\subsection{First- and Second-Order Conditions}

\begin{theorem}
\label{thm:optimality}
$f$가 $x_\ast$의 근방에서 두 번 연속 미분 가능하다고 하자.
\begin{enumerate}
\item $x_\ast$가 국소 최소점이면 $\nabla f(x_\ast)=0$이고 $\nabla^2f(x_\ast)$는 양반정치이다.
\item $\nabla f(x_\ast)=0$이고 $\nabla^2f(x_\ast)$가 양정치이면 $x_\ast$는 엄격한 국소 최소점이다.
\end{enumerate}
\end{theorem}

\begin{proof}
(i) 방향 $d$를 하나 잡고 $\phi(t)=f(x_\ast+td)$로 두면 $\phi$는 $t=0$에서 국소 최소이므로 $\phi'(0)=\nabla f(x_\ast)^{\mathsf T}d=0$이고 $d$가 임의이므로 $\nabla f(x_\ast)=0$이다. Taylor 정리에서 어떤 $\theta\in(0,1)$에 대해 $\phi(t)=\phi(0)+(t^2/2)\,d^{\mathsf T}\nabla^2f(x_\ast+\theta td)d$이고 작은 $t$에서 $\phi(t)\ge\phi(0)$이므로 $d^{\mathsf T}\nabla^2f(x_\ast+\theta td)d\ge0$이다. $t\to0$으로 보내면 연속성에서 $d^{\mathsf T}\nabla^2f(x_\ast)d\ge0$이다.

(ii) $\nabla^2f(x_\ast)$의 가장 작은 고유값을 $\lambda>0$이라 하자. 고유값은 행렬 원소의 연속 함수이므로(\cref{cor:weyl}) $x_\ast$의 어떤 근방에서 $\nabla^2f$의 가장 작은 고유값이 $\lambda/2$ 이상이다. 그 근방의 $x\ne x_\ast$에서 Taylor 정리로 $f(x)=f(x_\ast)+(1/2)(x-x_\ast)^{\mathsf T}\nabla^2f(\xi)(x-x_\ast)\ge f(x_\ast)+(\lambda/4)\|x-x_\ast\|^2>f(x_\ast)$이다.
\end{proof}

$\nabla f(x)=0$인 점을 \term{정류점}{stationary point}이라 한다. 정류점은 최소점일 수도, 최대점일 수도, 안장점일 수도 있다(\cref{fig:stationary}). 두 조건 사이에는 틈이 있다. $x^4$는 $0$에서 최소이지만 이계 도함수가 $0$이고 $x^3$은 같은 조건을 만족하면서 최소가 아니다. 기울기만 보는 알고리즘은 정류점을 찾을 뿐이므로 찾은 점이 최소점인지는 함수의 구조(뒤에서 볼 볼록성)나 이계 조건에 기대어 판단한다.

\begin{figure}[htbp]
\centering
\begin{tikzpicture}
\begin{axis}[sdplot,width=0.42\textwidth,height=0.42\textwidth,xlabel={$x_1$},ylabel={$x_2$},xmin=-2,xmax=2,ymin=-2,ymax=2,axis equal image,title={최소점: $x_1^2+3x_2^2$}]
\foreach \c in {0.5,1.5,3,5,8} {\addplot[black,domain=0:360,samples=80,variable=\t] ({sqrt(\c)*cos(\t)},{sqrt(\c/3)*sin(\t)});}
\addplot[only marks,mark=*,mark size=1.5pt] coordinates {(0,0)};
\end{axis}
\end{tikzpicture}\hfill
\begin{tikzpicture}
\begin{axis}[sdplot,width=0.42\textwidth,height=0.42\textwidth,xlabel={$x_1$},ylabel={$x_2$},xmin=-2,xmax=2,ymin=-2,ymax=2,axis equal image,title={안장점: $x_1^2-x_2^2$}]
\foreach \c in {0.5,1.5,3} {\addplot[black,domain=-1.4:1.4,samples=60,variable=\s] ({sqrt(\c)*cosh(\s)},{sqrt(\c)*sinh(\s)}); \addplot[black,domain=-1.4:1.4,samples=60,variable=\s] ({-sqrt(\c)*cosh(\s)},{sqrt(\c)*sinh(\s)});
\addplot[black,dashed,domain=-1.4:1.4,samples=60,variable=\s] ({sqrt(\c)*sinh(\s)},{sqrt(\c)*cosh(\s)}); \addplot[black,dashed,domain=-1.4:1.4,samples=60,variable=\s] ({sqrt(\c)*sinh(\s)},{-sqrt(\c)*cosh(\s)});}
\addplot[only marks,mark=*,mark size=1.5pt] coordinates {(0,0)};
\end{axis}
\end{tikzpicture}
\caption{두 정류점의 등고선. 왼쪽은 Hessian이 양정치인 최소점, 오른쪽은 Hessian의 고유값이 $2$와 $-2$인 안장점이다(실선은 양의 값, 점선은 음의 값의 등고선). 둘 다 기울기가 $0$이다.}
\label{fig:stationary}
\end{figure}

\subsection{Stopping Rules as Optimality Tests}

반복법은 정확한 정류점에 닿지 않으므로 ``기울기가 충분히 작다''를 수로 정해야 한다. stardust의 \code{minimize}는 방법마다 scipy와 같은 판정을 쓴다. 결과의 \code{converged}는 그 판정을 통과했다는 표시일 뿐 수학적인 최적성을 보증하지 않는다. BFGS는 $\|\nabla f(x)\|_\infty\le\mathrm{gtol}$(기본 $10^{-5}$)이고 L-BFGS-B는 같은 자리에 \cref{sec:kkt}의 사영 기울기를 쓴다. 이 판정에는 척도가 없다. $f$에 $10^6$을 곱하면 같은 점의 기울기도 $10^6$배가 되므로 절대 문턱은 문제의 단위에 따라 너무 엄하거나 너무 느슨하다. 강볼록이면 \cref{sec:strong-convexity}에서 보듯이 $\|\nabla f\|$가 함수값의 차와 거리를 모두 위에서 막아 주므로 문턱의 뜻이 분명해진다. 일반적인 함수에서는 그런 보증이 없어 사용자가 문제의 척도에 맞게 \code{gtol}을 고르거나 결과를 다른 방법으로 확인해야 한다. 이 장의 뒷부분에서 볼 쌍대 간격은 척도가 있는 판정의 예이다.

멈춤 판정에는 기울기 말고 다른 것도 있다. stardust가 방법마다 쓰는 판정을 \cref{tab:stop-rules}에 모았다. L-BFGS-B는 사영 기울기 판정 말고도 함수값의 상대 감소 $f_k-f_{k+1}\le\mathrm{ftol}\cdot\max(|f_k|,|f_{k+1}|,1)$(기본 $\mathrm{ftol}=2.2\times10^{-9}$)이면 멈추고 이 판정이 대개 먼저 걸린다. \cref{tab:box-kkt}의 상자 문제를 기본값으로 풀면 \code{gtol}을 $10^{-11}$까지 줄여도 $19$번째 반복에서 사영 기울기 $1.0\times10^{-3}$으로 멈추고 \code{converged}는 참이다. 상대 감소 판정을 통과했지만 그 반복이 실제로 걸음을 내딛지 못한 경우(선 탐색이 앞의 점으로 돌아온 경우)를 scipy는 수렴으로 보고하고 stardust는 사영 기울기가 문턱 위인 한 \code{converged}를 거짓으로 둔다. \code{converged}에는 함수값이 유한하다는 조건도 붙는다.

\begin{table}[htbp]
\begin{sdtable}{colspec={X[1.5,l]X[3,l]X[1.6,l]}, row{1}={font=\small\bfseries}, cells={font=\small}}
\toprule
방법 & 멈춤 판정 & 기본값 \\
\midrule
Nelder--Mead & 단체의 지름과 꼭짓점 함수값의 폭 & $10^{-4}$, $10^{-4}$ \\
BFGS & $\|\nabla f\|_\infty\le\mathrm{gtol}$ 또는 걸음 $\alpha\|p\|\le\mathrm{xrtol}(\mathrm{xrtol}+\|x\|)$ & $10^{-5}$, $0$ \\
L-BFGS-B & 사영 기울기 $\le\mathrm{gtol}$ 또는 상대 감소 $\le\mathrm{ftol}$(걸음이 없으면 수렴 아님) & $10^{-5}$, $2.2\times10^{-9}$ \\
Newton-CG & 걸음의 $\|\cdot\|_1\le n\cdot\mathrm{xtol}$ & $10^{-5}$ \\
trust-ncg, trust-exact & $\|\nabla f\|<\mathrm{gtol}$ & $10^{-4}$ \\
\code{least\_squares} & 비용·걸음·기울기의 상대 판정 & $10^{-8}$ 셋 \\
\code{lsq\_linear} & KKT 잔차(방법마다 척도가 다름) 또는 비용의 상대 변화 & $10^{-10}$ \\
\code{nnls} & \cref{eq:nnls-kkt}를 정확히, 반복 상한 & $3n-1$ \\
\code{minimize\_scalar} & 구간 폭(\code{'brent'} $1.48\times10^{-8}$, \code{'bounded'} $\sqrt\epsilon|x|+\mathrm{xatol}/3$) & \cref{sec:opt-precision} \\
\code{regression.lasso} & 쌍대 간격 $\le\mathrm{tol}\cdot\|v\|_\Omega^2$ & $10^{-4}$ \\
\bottomrule
\end{sdtable}
\caption{stardust의 멈춤 판정(코드와 설계 문서에서 확인; 9권에서 방법마다 자세히 다룬다). 척도가 있는 판정은 \code{lasso}의 간격과 \code{nnls}의 정확한 KKT뿐이다.}
\label{tab:stop-rules}
\end{table}

\section{Convex Sets}
\label{sec:convex-sets}

\subsection{Definitions and Operations}

집합 $C\subseteq\mathbb R^n$이 그 안의 두 점을 잇는 선분을 모두 포함하면, 곧 $x,y\in C$와 $t\in[0,1]$에서 $tx+(1-t)y\in C$이면 \term{볼록 집합}{convex set}이라 한다. 점들의 \term{볼록 결합}{convex combination} $\sum_i\lambda_ix_i$($\lambda_i\ge0$, $\sum_i\lambda_i=1$)을 모두 모은 것이 그 점들의 \term{볼록 껍질}{convex hull}이다. 볼록 집합은 자기 점들의 볼록 결합에 닫혀 있다(점의 수에 대한 귀납법). 반공간 $\{x:a^{\mathsf T}x\le b\}$, 공, 상자 $\{l\le x\le u\}$, 양반정치 행렬의 집합은 볼록이다. 볼록 집합의 교집합, 아핀 사상의 상과 역상도 볼록이어서 선형 부등식 여러 개로 정한 다면체 $\{x:Ax\le b\}$는 볼록이다.

\begin{proposition}[Carath\'eodory]
\label{prop:caratheodory}
$S\subseteq\mathbb R^n$의 볼록 껍질에 든 점은 $S$의 점 $n+1$개 이하의 볼록 결합이다.
\end{proposition}

\begin{proof}
$x=\sum_{i=1}^k\lambda_ix_i$($\lambda_i>0$)이고 $k>n+1$이라 하자. 벡터 $x_2-x_1,\dots,x_k-x_1$은 $k-1>n$개이므로 일차종속이고 $\sum_i\mu_ix_i=0$, $\sum_i\mu_i=0$인, 모두 $0$은 아닌 $\mu$가 있다. 합이 $0$이므로 양수인 $\mu_i$가 있고 $\tau=\min\{\lambda_i/\mu_i:\mu_i>0\}$로 두면 $\lambda_i-\tau\mu_i$는 모두 음이 아니고 적어도 하나가 $0$이며 합은 여전히 $1$이다. 따라서 $x=\sum_i(\lambda_i-\tau\mu_i)x_i$는 점 하나를 덜 쓴 볼록 결합이고 이것을 $k\le n+1$이 될 때까지 되풀이한다.
\end{proof}

\subsection{Projection onto Closed Convex Sets}

제약이 있는 최적화에서 가장 자주 쓰는 연산은 집합 밖의 점을 집합 안의 가장 가까운 점으로 옮기는 사영이다.

\begin{theorem}[사영 정리]
\label{thm:projection}
$C\subseteq\mathbb R^n$이 공집합이 아닌 닫힌 볼록 집합이면 모든 $x$에 대해 $\|x-p\|$를 최소로 하는 $p\in C$가 하나뿐이다. 이 점을 $P_C(x)$로 쓰면 $p=P_C(x)$일 필요충분조건은 모든 $y\in C$에서 $(x-p)^{\mathsf T}(y-p)\le0$인 것이다. 또 $P_C$는 $\|P_C(x)-P_C(z)\|\le\|x-z\|$를 만족한다.
\end{theorem}

\begin{proof}
존재: $c_0\in C$를 하나 잡으면 $\{y\in C:\|x-y\|\le\|x-c_0\|\}$은 공집합이 아닌 컴팩트 집합이고 그 위에서 연속 함수 $\|x-y\|$가 최솟값을 갖는다. 특징: $p$가 최소점이고 $y\in C$이면 $t\in[0,1]$에서 $p+t(y-p)\in C$이므로 $\|x-p-t(y-p)\|^2\ge\|x-p\|^2$, 곧 $-2t(x-p)^{\mathsf T}(y-p)+t^2\|y-p\|^2\ge0$이고 $t\to0^+$에서 부등식을 얻는다. 거꾸로 부등식이 성립하면 $\|x-y\|^2=\|x-p\|^2-2(x-p)^{\mathsf T}(y-p)+\|y-p\|^2\ge\|x-p\|^2$이다. 유일성: $p$와 $q$가 모두 최소점이면 두 부등식 $(x-p)^{\mathsf T}(q-p)\le0$과 $(z-q)^{\mathsf T}(p-q)\le0$($z=x$)을 더해 $\|p-q\|^2\le0$이다. 비확대성: $p=P_C(x)$, $q=P_C(z)$에 대해 $(x-p)^{\mathsf T}(q-p)\le0$과 $(z-q)^{\mathsf T}(p-q)\le0$을 더하면 $\|p-q\|^2\le(x-z)^{\mathsf T}(p-q)\le\|x-z\|\|p-q\|$이다.
\end{proof}

부등식 $(x-p)^{\mathsf T}(y-p)\le0$은 $x-p$가 $p$에서 $C$의 어느 점 $y$로 가는 방향 $y-p$와도 둔각(직각 포함)을 이룬다는 뜻이다(\cref{fig:projection}). 몇 가지 집합은 사영이 닫힌 꼴이다.
\begin{enumerate}
\item 상자 $\{l\le x\le u\}$: 성분마다 자른다, $P(x)_i=\min(u_i,\max(l_i,x_i))$. 제약이 성분마다 따로이므로 문제가 좌표별로 갈라진다.
\item 반공간 $\{a^{\mathsf T}x\le b\}$: $x$가 밖에 있으면 $x-((a^{\mathsf T}x-b)/\|a\|^2)\,a$이다.
\item 공 $\{\|x-c\|\le r\}$: 밖에 있으면 $c+r(x-c)/\|x-c\|$이다.
\end{enumerate}
상자 사영은 L-BFGS-B(\cref{sec:smooth}, 97장)와 \code{lsq\_linear}의 \code{'bvls'}(\cref{sec:nnls}, 101장)가 반복마다 쓴다(\code{'trf'}는 사영하지 않고 경계에서 반사하거나 내부로 밀어 넣는다). 성분마다 비교를 두 번 하므로 비용이 $O(n)$이고 반올림이 없다. L-BFGS-B는 성분마다 경계의 종류(없음, 아래만, 둘 다, 위만)를 정수 \code{nbd}로 따로 둔다. 없는 경계의 자리에는 $\pm\infty$ 대신 $0$이 들어 있어서 \code{nbd}를 보지 않으면 비교를 할 수 없다. 시작점이 상자 밖이면 첫 반복 전에 사영한다.

\subsection{Separation and Farkas' Lemma}

볼록 해석의 기본 도구인 분리 정리는 사영 정리에서 나온다.

\begin{theorem}[분리 정리]
\label{thm:separation}
$C$가 공집합이 아닌 닫힌 볼록 집합이고 $x\notin C$이면 $a^{\mathsf T}x>\sup_{y\in C}a^{\mathsf T}y$인 $a$가 있다. $C$가 볼록이고(닫혀 있지 않아도 된다) $x$가 $C$의 내부에 있지 않으면 $a\ne0$이고 $a^{\mathsf T}x\ge\sup_{y\in C}a^{\mathsf T}y$인 $a$가 있다(지지 초평면).
\end{theorem}

\begin{proof}
첫째 주장: $p=P_C(x)$, $a=x-p\ne0$으로 두면 \cref{thm:projection}에서 모든 $y\in C$에 대해 $a^{\mathsf T}y\le a^{\mathsf T}p=a^{\mathsf T}x-\|a\|^2$이다. 둘째 주장: 먼저 $x$가 $\bar C$의 내부에도 있지 않음을 보인다. $C$의 내부가 비어 있으면 $C$의 아핀 껍질이 진 아핀 부분 공간이므로 그 법선 $a$를 부호를 맞춰 잡으면($C$ 위에서 $a^{\mathsf T}y$는 상수이다) $a^{\mathsf T}x\ge a^{\mathsf T}y$가 되어 바로 지지 초평면을 준다. 내부의 점 $c_0$가 있으면 $y\in\bar C$일 때 반열린 선분 $[c_0,y)$가 $C$의 내부에 든다는 선분 원리를 쓴다($c_0$ 둘레의 공을 $y$ 쪽으로 축소해 옮긴 것들이 $C$에 들기 때문이다). $x$가 $\bar C$의 내부에 있다면 작은 $\epsilon>0$에서 $y=x+\epsilon(x-c_0)\in\bar C$이고 $x$는 $[c_0,y)$ 위에 있으므로 $C$의 내부에 있게 되어 가정에 어긋난다(Rockafellar\cite{rockafellar1970} 정리 6.1, 6.3). 그러므로 $\bar C$ 밖에서 $x$로 수렴하는 점열 $x_k$가 있다. 각 $x_k$를 $\bar C$에서 분리하는 단위 벡터 $a_k$를 잡으면 부분수열이 단위 벡터 $a$로 수렴하고 극한에서 $a^{\mathsf T}x\ge a^{\mathsf T}y$($y\in C$)이다.
\end{proof}

\begin{lemma}[Farkas]
\label{lem:farkas}
행렬 $A\in\mathbb R^{m\times n}$과 $b\in\mathbb R^m$에 대해 다음 둘 가운데 정확히 하나가 성립한다. (가) $Ax=b$, $x\ge0$인 $x$가 있다. (나) $A^{\mathsf T}y\le0$이고 $b^{\mathsf T}y>0$인 $y$가 있다.
\end{lemma}

\begin{proof}
둘이 함께 성립하면 $0<b^{\mathsf T}y=x^{\mathsf T}A^{\mathsf T}y\le0$이므로 모순이다. (가)가 성립하지 않으면 $b$는 원뿔 $K=\{Ax:x\ge0\}$ 밖에 있다. $K$는 유한개의 벡터가 생성한 원뿔이므로 닫혀 있다(\cref{prop:caratheodory}의 논법을 원뿔에 쓰면 $K$는 일차독립인 열들이 생성한 원뿔들의 유한 합집합이고 열이 독립인 $A_S$에서는 $\|A_Sy\|\ge\sigma_{\min}(A_S)\|y\|$이므로 $\{A_Sy:y\ge0\}$이 닫혀 있다). \cref{thm:separation}에서 $y^{\mathsf T}b>\sup_{z\in K}y^{\mathsf T}z$인 $y$가 있다. $0\in K$이므로 $y^{\mathsf T}b>0$이고 $K$가 원뿔이므로 상한은 $0$이어야 하며 따라서 모든 열 $a_j$에서 $y^{\mathsf T}a_j\le0$, 곧 $A^{\mathsf T}y\le0$이다.
\end{proof}

Farkas가 1902년에 증명한 이 정리\cite{farkas1902}는 ``해가 없다''를 증명서 $y$ 하나로 보여 준다. 이 장의 쌍대 이론은 모두 이 꼴을 따른다. 해가 있음은 해를 내밀어 보이면 되지만 해가 없음은 그 자체로는 보일 방법이 없다. Farkas는 그 부재를 다른 체계의 해가 있음으로 바꾸어 놓았다. 선형 계획의 쌍대 정리와 심플렉스법, 내부점법은 103장에서 이 보조정리 위에 세운다.

\begin{figure}[htbp]
\centering
\begin{tikzpicture}[x=9mm,y=9mm,every node/.style={font=\small}]
\fill[black!12] (0,0) -- (3,0) -- (4,1.6) -- (2.6,3) -- (0.4,2.4) -- cycle;
\draw[thick] (0,0) -- (3,0) -- (4,1.6) -- (2.6,3) -- (0.4,2.4) -- cycle;
\node at (1.6,1.2) {$C$};
\coordinate (x) at (4.9,3.9);
\coordinate (p) at (3.3,2.3);
\fill (x) circle (1.6pt) node[above right] {$x$};
\fill (p) circle (1.6pt) node[left=3pt] {$P_C(x)$};
\draw[->,thick] (p) -- (x);
\draw[dashed] (2.2,3.4) -- (5.0,0.6);
\node[anchor=north west] at (5.0,0.6) {분리 초평면};
\fill (1.2,0.6) circle (1.2pt) node[below] {$y$};
\draw[->] (p) -- (1.2,0.6);
\end{tikzpicture}
\caption{닫힌 볼록 집합으로의 사영. $x-P_C(x)$는 $P_C(x)$에서 $C$의 어느 점 $y$로 가는 방향과도 직각 이상의 각을 이루고(\cref{thm:projection}) $P_C(x)$를 지나며 $x-P_C(x)$에 수직인 초평면이 $x$와 $C$를 분리한다(\cref{thm:separation}).}
\label{fig:projection}
\end{figure}

\subsection{Cones and Normal Cones}

$t\ge0$을 곱해도 닫혀 있는 집합을 \term{원뿔}{cone}이라 한다. 볼록 집합 $C$와 그 점 $x$에서
\[
N_C(x)=\{v:v^{\mathsf T}(y-x)\le0\ \text{for all }y\in C\}
\]
를 $x$에서의 \term{법선 원뿔}{normal cone}이라 한다. 내부의 점에서는 $\{0\}$이고 경계에서는 바깥쪽을 향한 벡터들이다. \cref{thm:projection}의 조건은 $x-P_C(x)\in N_C(P_C(x))$로 다시 쓸 수 있다. 상자에서는 계산이 성분별로 갈라진다. $l_i<x_i<u_i$이면 $v_i=0$, $x_i=l_i$이면 $v_i\le0$, $x_i=u_i$이면 $v_i\ge0$이 $v\in N_C(x)$의 조건이다(\cref{fig:normal-cone}). 제약이 있을 때 ``기울기가 $0$''을 대신하는 조건 $-\nabla f(x)\in N_C(x)$(\cref{sec:kkt})는 법선 원뿔로 적는다.

\section{Convex Functions}
\label{sec:convex-functions}

\subsection{Definitions and First-Order Characterization}

볼록 집합 $C$ 위의 함수 $f$가
\[
f\bigl(tx+(1-t)y\bigr)\le tf(x)+(1-t)f(y)\qquad(x,y\in C,\ t\in[0,1])
\]
를 만족하면 \term{볼록 함수}{convex function}라 한다. 그래프 위의 두 점을 이은 현이 그래프보다 아래로 내려가지 않는다. $x\ne y$, $t\in(0,1)$에서 부등식이 엄격하면 \term{엄격한 볼록}{strictly convex}이다. 같은 조건을 집합으로 말하면 \term{에피그래프}{epigraph} $\{(x,s):f(x)\le s\}$가 볼록 집합이다. 이 관점에서 $+\infty$ 값을 허용하면 제약을 함수에 넣을 수 있다. $C$ 위에서 $0$이고 밖에서 $+\infty$인 \term{지시 함수}{indicator function} $\iota_C$는 $C$가 볼록일 때 볼록이고 ``$C$ 위에서 $f$를 최소화''는 ``$f+\iota_C$를 최소화''와 같다. 볼록 결합으로 귀납하면 Jensen의 부등식 $f(\sum\lambda_ix_i)\le\sum\lambda_if(x_i)$를 얻고 확률로 쓰면 $f(EX)\le Ef(X)$이다(9장에서 측도론의 판으로 다시 본다).

\begin{theorem}
\label{thm:convex-characterization}
$f$가 열린 볼록 집합 위에서 미분 가능하면 다음은 동치이다.
\begin{enumerate}
\item $f$는 볼록이다.
\item 모든 $x,y$에서 $f(y)\ge f(x)+\nabla f(x)^{\mathsf T}(y-x)$이다.
\item 모든 $x,y$에서 $(\nabla f(y)-\nabla f(x))^{\mathsf T}(y-x)\ge0$이다(기울기의 단조성).
\end{enumerate}
$f$가 두 번 연속 미분 가능하면 이것은 다시 모든 $x$에서 $\nabla^2f(x)$가 양반정치인 것과 동치이다.
\end{theorem}

\begin{proof}
(i)$\Rightarrow$(ii): 볼록성에서 $f(x+t(y-x))\le f(x)+t(f(y)-f(x))$이므로 $(f(x+t(y-x))-f(x))/t\le f(y)-f(x)$이고 $t\to0^+$에서 좌변이 $\nabla f(x)^{\mathsf T}(y-x)$이다. (ii)$\Rightarrow$(i): $z=tx+(1-t)y$에서 (ii)를 $x$와 $y$에 각각 쓰면 $f(x)\ge f(z)+\nabla f(z)^{\mathsf T}(x-z)$, $f(y)\ge f(z)+\nabla f(z)^{\mathsf T}(y-z)$이고 두 식에 $t$와 $1-t$를 곱해 더하면 $t(x-z)+(1-t)(y-z)=0$이므로 $tf(x)+(1-t)f(y)\ge f(z)$이다. (ii)$\Rightarrow$(iii): (ii)를 $(x,y)$와 $(y,x)$에 쓰고 더한다. (iii)$\Rightarrow$(ii): $\phi(t)=f(x+t(y-x))$이면 $\phi'(t)-\phi'(0)=(\nabla f(x+t(y-x))-\nabla f(x))^{\mathsf T}(y-x)$이고 (iii)을 $x$와 $x+t(y-x)$에 쓰면 이것이 $t>0$에서 음이 아니므로 평균값 정리에서 어떤 $\theta\in(0,1)$에 대해 $\phi(1)-\phi(0)=\phi'(\theta)\ge\phi'(0)$이다. 이계 조건: $\nabla^2f\succeq0$이면 $\phi''(t)=(y-x)^{\mathsf T}\nabla^2f(x+t(y-x))(y-x)\ge0$이라 $\phi'$가 증가하고 (iii)이 나온다. 거꾸로 (iii)에서 $y=x+td$로 두고 $t^2$으로 나눈 뒤 $t\to0$으로 보내면 $d^{\mathsf T}\nabla^2f(x)d\ge0$이다.
\end{proof}

(ii)에 따르면 볼록 함수의 그래프는 어느 접평면보다도 아래로 내려가지 않는다(\cref{fig:jensen}). 그래서 정류점 $\nabla f(x_\ast)=0$에서 $f(y)\ge f(x_\ast)$가 모든 $y$에 대해 성립한다.

\begin{corollary}
\label{cor:convex-global}
볼록 함수의 국소 최소점은 전역 최소점이고 미분 가능하면 정류점과 전역 최소점은 같다. 최소점들의 집합은 볼록이고 엄격한 볼록 함수의 최소점은 많아야 하나이다.
\end{corollary}

\begin{proof}
$x_\ast$가 국소 최소점이고 $f(y)<f(x_\ast)$인 $y$가 있다면 선분 위의 점 $x_\ast+t(y-x_\ast)$에서 $f\le(1-t)f(x_\ast)+tf(y)<f(x_\ast)$이고 $t\to0$으로 보내면 국소 최소성에 어긋난다. 정류점에 관한 주장은 (ii)에서 나온다. 최소점 $x,y$와 최솟값 $f_\ast$에 대해 $f(tx+(1-t)y)\le f_\ast$이므로 선분이 최소점 집합에 들어가고 엄격하면 $f(tx+(1-t)y)<f_\ast$라 모순이 생긴다.
\end{proof}

\begin{figure}[htbp]
\centering
\begin{tikzpicture}
\begin{axis}[sdplot,xlabel={$x$},ylabel={$f(x)$},xmin=-2.2,xmax=2.6,ymin=-0.5,ymax=6]
\addplot[black,thick,domain=-2.2:2.6,samples=120] {ln(exp(x)+exp(-x/2))+0.3*x^2};
\addlegendentry{$f$}
\addplot[black,dashed,domain=-1.6:2] {(ln(exp(-1.6)+exp(0.8))+0.3*2.56)+((ln(exp(2)+exp(-1))+1.2)-(ln(exp(-1.6)+exp(0.8))+0.3*2.56))/3.6*(x+1.6)};
\addlegendentry{현}
\addplot[black!55,domain=-2.2:2.2] {(ln(exp(0.3)+exp(-0.15))+0.027)+((exp(0.3)-0.5*exp(-0.15))/(exp(0.3)+exp(-0.15))+0.18)*(x-0.3)};
\addlegendentry{접선}
\addplot[only marks,mark=*,mark size=1.3pt,samples at={-1.6,0.3,2}] {ln(exp(x)+exp(-x/2))+0.3*x^2};
\end{axis}
\end{tikzpicture}
\caption{볼록 함수 $f(x)=\log(e^x+e^{-x/2})+0.3x^2$의 현과 접선. 현은 그래프 위에(정의), 접선은 그래프 아래에(\cref{thm:convex-characterization}의 (ii)) 있다.}
\label{fig:jensen}
\end{figure}

\subsection{Operations That Preserve Convexity}

새로운 함수가 볼록인지 정의로 확인하는 일은 번거롭다. 다음 연산들이 볼록성을 지키므로 대부분의 함수는 아는 조각에서 조립해 판단한다. 음이 아닌 계수의 합, 아핀 사상과의 합성 $f(Ax+b)$, 볼록 함수들의 점별 상한 $\sup_\alpha f_\alpha$(에피그래프의 교집합), 증가하는 볼록 함수와 볼록 함수의 합성, 그리고 볼록 함수 $g(x,y)$를 $y$에 대해 볼록 집합 위에서 최소화한 $\inf_yg(x,y)$(에피그래프의 사영)이 모두 볼록이다. 이 장과 뒤의 권에서 만나는 목적 함수를 이렇게 확인할 수 있다.
\begin{enumerate}
\item 최소제곱 $(1/2)\|Ax-b\|^2$: Hessian $A^{\mathsf T}A\succeq0$.
\item 노름 $\|x\|$: 삼각부등식과 동차성에서 바로 볼록이다. $\|x\|_1$은 엄격한 볼록이 아니다.
\item log-sum-exp $\operatorname{lse}(x)=\log\sum_ie^{x_i}$: Hessian이 $\operatorname{diag}(\pi)-\pi\pi^{\mathsf T}$($\pi=\operatorname{softmax}(x)$)이고 $v^{\mathsf T}(\operatorname{diag}(\pi)-\pi\pi^{\mathsf T})v=\sum_i\pi_iv_i^2-(\sum_i\pi_iv_i)^2$는 확률 $\pi$ 아래에서 $v$의 분산이므로 음이 아니다.
\item 로지스틱 회귀의 음의 로그 우도 $\sum_i(\log(1+e^{x_i^{\mathsf T}\beta})-y_ix_i^{\mathsf T}\beta)$: 각 항이 $\operatorname{lse}(0,x_i^{\mathsf T}\beta)$와 일차식의 합이다.
\item 지수족 분포의 음의 로그 우도 $A(\eta)-\eta^{\mathsf T}T(x)$: 로그 분할 함수 $A(\eta)=\log\int e^{\eta^{\mathsf T}T(x)}h(x)\,dx$를 적분 아래에서 두 번 미분하면 $\nabla^2A(\eta)=\operatorname{Cov}_\eta(T)\succeq0$이므로 볼록이다(10권). log-sum-exp는 범주 분포의 로그 분할 함수이고 위 셋째 항목의 Hessian이 바로 그 공분산이다.
\end{enumerate}
밀도 $p$에 대해 $-\log p$가 볼록이면 $p$를 \term{로그 오목}{log-concave}이라 한다. 정규, 지수, 로지스틱, 균등 분포의 밀도가 그렇다. 로그 오목 밀도는 단봉이며 꼬리가 지수 함수보다 느리게 줄지 않는다. 위 준위 집합 $\{p\ge t\}=\{-\log p\le-\log t\}$가 볼록이므로 최빈값들의 집합은 볼록(한 점이거나 구간)이다. 또 적분 가능하면 $\phi=-\log p$는 강제적이어야 하는데 이런 볼록 함수는 어떤 $a>0$, $c$에 대해 $\phi(x)\ge a\|x\|-c$를 만족하므로 $p(x)\le e^{c-a\|x\|}$이다. 105장의 기각 표집과 12권의 모형 적합이 이 성질에 기댄다.

log-sum-exp는 볼록성과 수치 안정성이 만나는 좋은 예이다. $\operatorname{lse}(x)=m+\log\sum_ie^{x_i-m}$($m=\max_ix_i$)로 계산하면 지수가 넘치지 않고 합의 가장 큰 항이 $1$이어서 로그의 인자가 $[1,n]$에 머문다. stardust의 \code{special.logsumexp}(89장)도 이 꼴로 계산한다(\cref{alg:lse}). 이렇게 옮길 수 있는 것은 $\operatorname{lse}(x+c\mathbf 1)=\operatorname{lse}(x)+c$라는 평행 이동 성질 덕분이다. 이 성질을 Hessian으로 옮기면 영공간이 $\mathbf 1$ 방향이라는 말이 되고 lse가 엄격한 볼록이 아닌 방향도 이 방향이다.

\begin{algorithm}[htbp]
\begin{algorithmic}[1]
\Require 실수 배열 $x_1,\dots,x_n$
\Ensure $\operatorname{lse}(x)$
\State $m\gets\max_ix_i$ (NaN은 그대로 전한다); \textbf{if} $m=\pm\infty$ \textbf{then} $m\gets0$
\For{각 $i$ (NEON 두 칸씩, $1024$개 조각)}
  \State $(h_i,\ell_i)\gets\mathrm{TwoSum}(x_i,-m)$ \Comment{$x_i-m=h_i+\ell_i$ 정확히}
  \State $e_i\gets\exp(h_i)$; \textbf{if} $h_i$ 유한 \textbf{and} $e_i\ne0$ \textbf{then} $e_i\gets e_i+e_i\ell_i$ \Comment{$e^{h+\ell}\approx e^h(1+\ell)$}
\EndFor
\State $s\gets$ 조각별 쌍별 합을 정해진 순서로 더한 $\sum_ie_i$
\State \Return $m+\log s$
\end{algorithmic}
\caption{최댓값을 옮긴 log-sum-exp (stardust \code{csrc/special/reduce.cpp})}
\label{alg:lse}
\end{algorithm}

3행의 TwoSum을 빼고 $x_i-m$을 그냥 반올림하면 그 오차 $|x_i-m|\cdot2^{-53}$이 지수에 그대로 들어가 $x_i-m=-700$에서 $e^{x_i-m}$의 상대오차가 $7.8\times10^{-14}$가 된다. 차를 정확히 두 조각으로 나누고 작은 조각을 일차 보정으로 넣으면 그 오차가 사라진다. 최댓값이 무한이면 scipy의 관례를 따라 $m=0$으로 두어 $\infty-\infty$의 NaN을 피한다. 조각의 크기와 합의 순서가 입력 길이로만 정해지므로 결과는 스레드 수와 무관하다. 오차는 절대오차로 $u$의 작은 배수이므로 값 자체가 $0$ 근처이면 상대 정확도를 기대할 수 없다. 따로 빌드한 stardust에서 $\operatorname{lse}(0,-40)$은 참값 $4.25\times10^{-18}$ 대신 $0$이고 $\operatorname{lse}(-10^{-20},-50)$은 상대오차가 $2$퍼센트였다(\code{figures/08/lse.py}, 십진 참값). $s=1+(\text{작은 수})$를 binary64로 만드는 순간 작은 수가 반올림으로 사라지기 때문이다. 이런 값이 필요하면 최댓값 항을 빼고 $m+\log1p(\sum_{i\ne k}e^{x_i-m})$로 계산해야 한다.

\subsection{Strong Convexity and Smoothness}
\label{sec:strong-convexity}

수렴 속도와 멈춤 판정의 뜻을 말하려면 곡률의 위아래를 묶는 두 상수가 필요하다.

\begin{definition}
\label{def:strong-smooth}
미분 가능한 $f$가 모든 $x,y$에서 $f(y)\ge f(x)+\nabla f(x)^{\mathsf T}(y-x)+(\mu/2)\|y-x\|^2$을 만족하면 $\mu$-\term{강볼록}{strongly convex}이라 하고 $\|\nabla f(x)-\nabla f(y)\|\le L\|x-y\|$이면 $L$-\term{매끄럽다}{$L$-smooth}고 한다. 비 $\kappa=L/\mu$를 문제의 \term{조건수}{condition number}라 한다.
\end{definition}

$f$가 두 번 미분 가능하면 강볼록성은 $\nabla^2f\succeq\mu I$와 같고 볼록이면서 $L$-매끄러운 것은 $\nabla^2f\preceq LI$와 같다(\cref{thm:convex-characterization}의 증명을 $f-\mu\|x\|^2/2$와 $L\|x\|^2/2-f$에 쓴다. $L$-매끄러움은 $\|\nabla^2f\|_2\le L$과 같으므로 $\nabla^2f\succeq0$이면 $\nabla^2f\preceq LI$와 같다). 이차 함수 $(1/2)x^{\mathsf T}Ax-b^{\mathsf T}x$에서 $\mu$와 $L$은 $A$의 가장 작은 고유값과 가장 큰 고유값이고 $\kappa$는 4장의 $\kappa_2(A)$(\cref{def:cond})이다. \cref{thm:courant-fischer}로 보면 $\mu$와 $L$은 Rayleigh 몫 $d^{\mathsf T}\nabla^2f\,d/\|d\|^2$의 아래·위 한계이기도 하다.

\begin{proposition}
\label{prop:descent-lemma}
$f$가 $L$-매끄러우면 모든 $x,y$에서
\[
f(y)\le f(x)+\nabla f(x)^{\mathsf T}(y-x)+\frac L2\|y-x\|^2
\]
이다(\term{하강 보조정리}{descent lemma}). $f$가 $\mu$-강볼록이고 최소점이 $x_\ast$이면
\[
\frac\mu2\|x-x_\ast\|^2\le f(x)-f(x_\ast)\le\frac1{2\mu}\|\nabla f(x)\|^2,\qquad\|x-x_\ast\|\le\frac1\mu\|\nabla f(x)\|
\]
이다.
\end{proposition}

\begin{proof}
첫 부등식: $f(y)-f(x)-\nabla f(x)^{\mathsf T}(y-x)=\int_0^1(\nabla f(x+t(y-x))-\nabla f(x))^{\mathsf T}(y-x)\,dt$이고 Cauchy--Schwarz와 Lipschitz 조건으로 피적분 함수는 $Lt\|y-x\|^2$ 이하이며 $\int_0^1t\,dt=1/2$이다. 강볼록의 첫 부등식은 정의에 $x\leftarrow x_\ast$, $\nabla f(x_\ast)=0$을 넣은 것이다. 둘째 부등식: 정의의 우변을 $y$에 대해 최소화하면 $y=x-\nabla f(x)/\mu$에서 값이 $f(x)-\|\nabla f(x)\|^2/(2\mu)$이므로 $f(x_\ast)\ge f(x)-\|\nabla f(x)\|^2/(2\mu)$이다. 셋째: 강볼록성의 정의를 $(x,x_\ast)$와 $(x_\ast,x)$에 쓰고 더하면 $(\nabla f(x)-\nabla f(x_\ast))^{\mathsf T}(x-x_\ast)\ge\mu\|x-x_\ast\|^2$이고 Cauchy--Schwarz를 쓴다.
\end{proof}

\cref{sec:optimality}에서 말한 ``척도가 있는 판정''은 둘째 줄에서 나온다. 강볼록 상수 $\mu$를 알면 $\|\nabla f\|\le\epsilon$은 거리 $\epsilon/\mu$와 함수값의 차 $\epsilon^2/(2\mu)$를 보증한다. 반대로 $\mu$가 작은 문제, 곧 조건수가 큰 문제에서는 작은 기울기도 먼 거리와 양립하므로 기울기 문턱만으로는 해의 정확도를 말할 수 없다.

\section{Subgradients and Proximal Steps}
\label{sec:subgradients}

\subsection{Subdifferentials}

라소의 $\|b\|_1$처럼 꺾인 곳이 있는 볼록 함수에서는 관심 있는 점이 하필 기울기가 정의되지 않는 점이다. \cref{thm:convex-characterization}의 (ii)가 기울기를 대신할 개념을 준다.

\begin{definition}
\label{def:subgradient}
볼록 함수 $f$와 그 정의역의 점 $x$에 대해 모든 $y$에서 $f(y)\ge f(x)+g^{\mathsf T}(y-x)$인 $g$를 $x$에서의 \term{부분기울기}{subgradient}라 하고 그 전체를 \term{부분미분}{subdifferential} $\partial f(x)$로 쓴다.
\end{definition}

부분미분은 닫힌 볼록 집합이고 $f$가 $x$에서 미분 가능하면 $\{\nabla f(x)\}$ 하나뿐이다. $g\in\partial f(x)$이면 정의에서 $(f(x+td)-f(x))/t\ge g^{\mathsf T}d$이고 $t\to0^+$에서 $\nabla f(x)^{\mathsf T}d\ge g^{\mathsf T}d$가 모든 $d$에서 성립하는데 $-d$를 넣으면 등호이므로 $g=\nabla f(x)$이다. 정의역의 내부에서는 공집합이 아니다. 에피그래프의 경계점 $(x,f(x))$에서 \cref{thm:separation}의 지지 초평면을 잡고 그 법선을 $(g,-1)$꼴로 맞추면 된다(법선의 마지막 성분이 $0$이면 초평면이 수직이 되어 $x$가 정의역의 내부라는 데 어긋난다). 정의를 그대로 읽으면 다음 최적성 조건이 된다.

\begin{proposition}
\label{prop:subgradient-optimality}
볼록 함수 $f$에 대해 $x_\ast$가 전역 최소점일 필요충분조건은 $0\in\partial f(x_\ast)$이다.
\end{proposition}

\begin{proof}
$0\in\partial f(x_\ast)$는 모든 $y$에서 $f(y)\ge f(x_\ast)$라는 말 그대로이다.
\end{proof}

절댓값의 부분미분은 $\partial|x|=\{\operatorname{sign}x\}$($x\ne0$), $\partial|0|=[-1,1]$이다(\cref{fig:soft}). $\|b\|_1=\sum_j|b_j|$은 좌표마다 따로인 함수의 합이므로 부분미분은 좌표별 부분미분의 곱집합이다. 두 볼록 함수의 합에서 $\partial f(x)+\partial g(x)\subseteq\partial(f+g)(x)$는 정의에서 나온다. 등호는 한쪽이 미분 가능하면 성립한다. $f$가 미분 가능할 때 $h\in\partial(f+g)(x)$이면 $y=x+t(z-x)$에서 $g(y)-g(x)\ge h^{\mathsf T}(y-x)-(f(y)-f(x))$이고 $g$의 볼록성 $g(y)\le(1-t)g(x)+tg(z)$와 합쳐 $t$로 나눈 뒤 $t\to0^+$으로 보내면 $g(z)-g(x)\ge(h-\nabla f(x))^{\mathsf T}(z-x)$, 곧 $h-\nabla f(x)\in\partial g(x)$이다. 일반적인 합의 규칙(Moreau--Rockafellar)은 두 정의역의 상대 내부가 만날 때 성립하며 분리 정리를 곱공간의 에피그래프에 써서 증명한다(Rockafellar\cite{rockafellar1970} 정리 23.8).

\subsection{The Lasso Optimality Conditions}

라소\cite{tibshirani1996}는 중심화하고 척도를 맞춘 설계 행렬 $Z$와 반응 $v$, 무게 $\Omega=\operatorname{diag}(\omega)$($\sum\omega_i=1$)에 대해
\begin{equation}
\label{eq:lasso}
\min_b\ \frac12\|v-Zb\|_\Omega^2+\lambda_1\|b\|_1+\frac{\lambda_2}2\|b\|^2,\qquad\|r\|_\Omega^2=\sum_i\omega_ir_i^2
\end{equation}
이다($\lambda_2>0$이면 엘라스틱넷). 첫 항이 미분 가능하므로 합의 규칙과 \cref{prop:subgradient-optimality}에서 $b$가 최소점일 필요충분조건은 잔차 $r=v-Zb$에 대해
\begin{equation}
\label{eq:lasso-kkt}
z_j^{\mathsf T}\Omega r-\lambda_2b_j\in\lambda_1\,\partial|b_j|\qquad(j=1,\dots,p)
\end{equation}
이다. 곧 $b_j\ne0$이면 $z_j^{\mathsf T}\Omega r-\lambda_2b_j=\lambda_1\operatorname{sign}b_j$이고 $b_j=0$이면 $|z_j^{\mathsf T}\Omega r|\le\lambda_1$이다. 잔차와 상관이 $\lambda_1$보다 작은 열은 계수가 정확히 $0$이 된다. 라소는 이렇게 변수를 고른다. 조건 \cref{eq:lasso-kkt}는 확인하기 쉬운 형태라서 stardust의 ``KKT 다듬기''(\cref{sec:duality})는 이 식을 그대로 판정으로 쓴다.

\subsection{Proximal Operators and Soft Thresholding}

꺾인 항이 좌표마다 따로이면 한 좌표에 대한 최소화를 닫힌 꼴로 풀 수 있다.

\begin{proposition}
\label{prop:soft}
$a>0$, $\lambda\ge0$일 때 $\phi(t)=(a/2)t^2-ct+\lambda|t|$의 최소점은
\[
t_\ast=\frac{S_\lambda(c)}a,\qquad S_\lambda(c)=\operatorname{sign}(c)\max(|c|-\lambda,0)
\]
이다. $S_\lambda$를 \term{연성 임계}{soft thresholding} 함수라 한다.
\end{proposition}

\begin{proof}
$\phi$는 강볼록이므로 최소점은 하나이고 \cref{prop:subgradient-optimality}에서 $0\in at_\ast-c+\lambda\,\partial|t_\ast|$이다. $t_\ast>0$이면 $at_\ast=c-\lambda$이므로 $c>\lambda$일 때 $t_\ast=(c-\lambda)/a$이고 $t_\ast<0$도 같은 식으로 $c<-\lambda$일 때이다. 남은 $|c|\le\lambda$에서는 $t_\ast=0$이 $c\in\lambda[-1,1]$로 조건을 만족한다.
\end{proof}

일반적으로 $\operatorname{prox}_g(x)=\arg\min_y\,g(y)+(1/2)\|y-x\|^2$를 $g$의 \term{근접 사상}{proximal operator}이라 하고 $g=\lambda|\cdot|$이면 $S_\lambda$, $g=\iota_C$이면 사영 $P_C$이다. 사영 정리는 근접 사상으로 그대로 넓어진다. 목적 함수가 강볼록이므로 최소점이 하나 있고 최적성 조건은 $p=\operatorname{prox}_g(x)\iff x-p\in\partial g(p)$이다. 부분미분은 단조이므로($a\in\partial g(p)$, $b\in\partial g(q)$이면 두 정의 부등식을 더해 $(a-b)^{\mathsf T}(p-q)\ge0$) $p=\operatorname{prox}_g(x)$, $q=\operatorname{prox}_g(z)$에 대해 $((x-p)-(z-q))^{\mathsf T}(p-q)\ge0$이고 항을 옮기면 $\|p-q\|^2\le(x-z)^{\mathsf T}(p-q)$이므로 Cauchy--Schwarz에서 비확대이다.

라소를 \term{좌표 하강}{coordinate descent}으로 풀면 좌표 하나의 최소화는 \cref{prop:soft}의 꼴이 된다\cite{friedman2010}. 다른 좌표를 고정하고 $b_j$를 $b_j+\delta$로 바꾸면 목적 함수는 $((d_j+\lambda_2)/2)(b_j+\delta)^2-(z_j^{\mathsf T}\Omega r+d_jb_j)(b_j+\delta)+\lambda_1|b_j+\delta|$에 상수를 더한 것이다($d_j=z_j^{\mathsf T}\Omega z_j$). 따라서 새 값은 $S_{\lambda_1}(z_j^{\mathsf T}\Omega r+d_jb_j)/(d_j+\lambda_2)$이고 잔차는 $r\gets r-z_j\delta$로 고친다(\cref{alg:cd}).

\begin{algorithm}[htbp]
\begin{algorithmic}[1]
\Require $Z$의 열 $z_j$, $d_j=z_j^{\mathsf T}\Omega z_j$, 잔차 $r=v-Zb$, $\lambda_1$, $\lambda_2$, 판정 여유 $\epsilon_j$
\Ensure 한 번 쓸고 난 $b$와 $r$, 가장 큰 변화 $\Delta_{\max}$
\State $\Delta_{\max}\gets0$
\For{$j=1,\dots,p$} \Comment{열 순서대로}
  \State $g\gets z_j^{\mathsf T}\Omega r$; $c\gets g+d_jb_j$ \Comment{공분산 갱신이면 $g=c_j-(Gb)_j$}
  \State \textbf{if} $|c|-\lambda_1>\epsilon_j$ \textbf{then} $b'\gets\operatorname{sign}(c)(|c|-\lambda_1)/(d_j+\lambda_2)$ \textbf{else} $b'\gets0$
  \If{$b'\ne b_j$}
    \State $r\gets r-z_j(b'-b_j)$; $\Delta_{\max}\gets\max(\Delta_{\max},|b'-b_j|)$; $b_j\gets b'$
  \EndIf
\EndFor
\end{algorithmic}
\caption{라소·엘라스틱넷의 좌표 하강 한 번 쓸기 (stardust \code{csrc/regression/penalized.cpp}의 \code{sweep})}
\label{alg:cd}
\end{algorithm}

\cref{alg:cd}의 각 걸음은 목적 함수를 늘리지 않으므로 함수값이 단조 감소한다. 또 목적 함수가 강제적이고 꺾인 항이 좌표마다 따로이므로 쓸기를 되풀이한 점열의 극한점은 모두 최소점이다(최소점이 하나이면 그 점으로 수렴한다; Tseng\cite{tseng2001}). 최소점이 하나인 강볼록 경우($\lambda_2>0$이거나 $Z$의 열이 독립)는 직접 보일 수 있다. 좌표 걸음마다 목적 함수 $F$가 $((d_j+\lambda_2)/2)\delta_j^2$ 이상 줄므로(좌표 하나의 함수가 그 곡률로 강볼록이다) $\sum\delta_j^2<\infty$이고 이웃한 반복점의 차가 $0$으로 간다. $F$가 강제적이므로 극한점이 있고 연성 임계가 연속이므로 극한점에서 모든 좌표가 \cref{prop:soft}의 조건을 만족한다. 꺾인 항이 좌표마다 따로이므로 $\partial F=\nabla f+\prod_j\partial(\lambda_1|b_j|)$이고 좌표마다의 조건이 곧 $0\in\partial F$이다. 최소점이 하나이므로 점열 전체가 그 점으로 수렴한다. 계수가 $0$인 열은 갱신에서 $g$ 하나만 계산하면 되므로 stardust는 0이 아닌 계수만 도는 짧은 쓸기를 그 계수들이 자리 잡을 때까지 되풀이하고 그 뒤에 전체를 한 번 쓴다. 열이 적으면($p\le1024$, \code{kCovMaxCols}) Gram 행렬 $G=Z^{\mathsf T}\Omega Z$를 한 번 만들어 두고 $g=c_j-(Gb)_j$로 계산한다. 그러면 걸음 하나의 비용이 행의 수 $m$에서 열의 수로 준다. 4행의 여유 $\epsilon_j$는 반올림을 감안한 문턱이다. $g$의 반올림 오차는 $u\sqrt{d_j\,v^{\mathsf T}\Omega v}$ 정도이므로 그보다 작은 차이로 문턱을 넘은 계수는 문턱 위에 있다고 말할 근거가 없다. stardust는 이 크기의 $64$배인 $\epsilon_j=64u\sqrt{d_j\,v^{\mathsf T}\Omega v}$를 여유로 두어 그런 열을 정확히 $0$에 묶어 둔다($\lambda_1=0$이면 여유도 $0$이고 상수 열 $d_j=0$은 건너뛴다). 활성 열을 그대로 복사한 열처럼 수학적으로 정확히 문턱 위에 놓이는 열은 이 여유가 없으면 반올림 때문에 $10^{-15}$짜리 계수를 얻는다.

\begin{figure}[htbp]
\centering
\begin{tikzpicture}
\begin{axis}[sdplot,width=0.44\textwidth,height=0.4\textwidth,xlabel={$x$},ylabel={$\partial|x|$},xmin=-2,xmax=2,ymin=-1.5,ymax=1.5]
\addplot[black,thick] coordinates {(-2,-1) (0,-1)};
\addplot[black,thick] coordinates {(0,1) (2,1)};
\addplot[black,very thick] coordinates {(0,-1) (0,1)};
\end{axis}
\end{tikzpicture}\hfill
\begin{tikzpicture}
\begin{axis}[sdplot,width=0.44\textwidth,height=0.4\textwidth,xlabel={$c$},ylabel={$S_\lambda(c)$},xmin=-2,xmax=2,ymin=-1.5,ymax=1.5]
\addplot[black,thick] coordinates {(-2,-1.4) (-0.6,0) (0.6,0) (2,1.4)};
\addplot[black!45,dashed,domain=-2:2] {x};
\node[font=\scriptsize,anchor=north] at (axis cs:0.6,0) {$\lambda$};
\node[font=\scriptsize,anchor=south] at (axis cs:-0.6,0) {$-\lambda$};
\end{axis}
\end{tikzpicture}
\caption{절댓값의 부분미분(왼쪽, $0$에서 구간 $[-1,1]$)과 연성 임계 함수 $S_\lambda$(오른쪽, $\lambda=0.6$; 점선은 항등 함수). $S_\lambda=(I+\lambda\,\partial|\cdot|)^{-1}$이다.}
\label{fig:soft}
\end{figure}

\subsection{Conjugate Functions}

쌍대 문제를 만드는 데 쓰는 도구가 켤레 함수이다. 함수 $f$의 \term{켤레}{conjugate}를
\[
f^\ast(y)=\sup_x\bigl(y^{\mathsf T}x-f(x)\bigr)
\]
로 정의한다. 일차 함수들의 상한이므로 $f$가 무엇이든 $f^\ast$는 볼록이고 정의에서 \term{Fenchel 부등식}{Fenchel's inequality} $f(x)+f^\ast(y)\ge x^{\mathsf T}y$가 나온다. 등호가 성립할 필요충분조건은 $y\in\partial f(x)$이다. $f(x)+f^\ast(y)=x^{\mathsf T}y$는 모든 $z$에서 $y^{\mathsf T}x-f(x)\ge y^{\mathsf T}z-f(z)$라는 말이고 항을 옮기면 그대로 부분기울기의 정의가 되기 때문이다. 이 장에서 쓰는 예는 셋이다. $f=(1/2)\|x\|^2$이면 $f^\ast=(1/2)\|y\|^2$이다. $f=\lambda\|x\|_1$이면 $f^\ast$는 $\|y\|_\infty\le\lambda$에서 $0$이고 밖에서 $+\infty$인 지시 함수이다($|y_j|>\lambda$인 좌표에서 $x_j$를 키우면 상한이 무한대이고 그렇지 않으면 $y^{\mathsf T}x\le\lambda\|x\|_1$이다). 그리고 $f=\iota_C$이면 $f^\ast(y)=\sup_{x\in C}y^{\mathsf T}x$로 $C$의 지지 함수이다. 닫힌 볼록 함수에서는 $f^{\ast\ast}=f$이다(Fenchel--Moreau 정리). 증명에서는 분리 정리를 써서 에피그래프를 그것을 담는 닫힌 반공간들의 교집합으로 나타낸다(Rockafellar\cite{rockafellar1970} 정리 12.2).

\section{Constrained Optimality: KKT Conditions}
\label{sec:kkt}

\subsection{Convex Sets as Constraints}

볼록 집합 $C$ 위에서 미분 가능한 볼록 함수 $f$를 최소화하는 문제의 최적성 조건은 \cref{prop:subgradient-optimality}를 $f+\iota_C$에 쓰면 얻는다. 정의에서 $\partial\iota_C(x)=N_C(x)$이고 합의 규칙에서 $\partial(f+\iota_C)(x)=\nabla f(x)+N_C(x)$이므로 $x_\ast$가 최소점일 필요충분조건은
\begin{equation}
\label{eq:normal-cone-condition}
-\nabla f(x_\ast)\in N_C(x_\ast)
\end{equation}
이다. 곧 모든 $y\in C$에서 $\nabla f(x_\ast)^{\mathsf T}(y-x_\ast)\ge0$이다. 내리막 방향은 모두 집합 밖을 향한다. \cref{thm:projection}과 견주면 이것은 $x_\ast=P_C(x_\ast-\alpha\nabla f(x_\ast))$($\alpha>0$ 아무것이나)와 같다. 사영 조건 $(x-p)^{\mathsf T}(y-p)\le0$에 $x=x_\ast-\alpha\nabla f(x_\ast)$, $p=x_\ast$를 넣으면 바로 \cref{eq:normal-cone-condition}이 되기 때문이다. 이 고정점 꼴이 계산에 쓰기 좋다.

\begin{definition}
\label{def:projected-gradient}
상자 $C=\{l\le x\le u\}$에서 $x$의 \term{사영 기울기}{projected gradient}를 $x-P_C(x-\nabla f(x))$로 정의한다. 그 무한 노름을 사영 기울기 노름이라 한다.
\end{definition}

사영 기울기가 $0$인 것은 \cref{eq:normal-cone-condition}과 같다. 성분으로 쓰면 $g=\nabla f(x)$에 대해 $g_i<0$(위로 가고 싶다)이면 위 경계까지의 거리로 자른 $\max(x_i-u_i,g_i)$, $g_i>0$이면 $\min(x_i-l_i,g_i)$이다. $x_i-g_i$가 상자 안에 있을 때에만 값이 $g_i$ 그대로이고 내부의 성분이라도 $|g_i|$가 경계까지의 거리보다 크면 잘린다. 경계에 붙어 있고 기울기가 밖을 향하는 성분은 $0$이 되어 판정에서 빠진다. 경계 쪽 성분의 크기가 경계까지의 거리로 잘리므로 이 노름은 ``얼마나 더 내려갈 수 있는지''를 경계를 넘지 않고 잰다.

\begin{algorithm}[htbp]
\begin{algorithmic}[1]
\Require $x$, $g=\nabla f(x)$, 경계 $l,u$와 경계의 종류 \code{nbd}$_i\in\{0,1,2,3\}$ (없음, 아래만, 둘 다, 위만)
\Ensure $\|x-P_C(x-g)\|_\infty$
\State $s\gets0$
\For{$i=1,\dots,n$}
  \State \textbf{if} $g_i$가 NaN \textbf{then} \Return $g_i$ \Comment{NaN이 판정을 통과하지 않게}
  \If{$g_i<0$ \textbf{and} 위 경계가 있다} $g_i\gets\max(x_i-u_i,\ g_i)$
  \ElsIf{$g_i\ge0$ \textbf{and} 아래 경계가 있다} $g_i\gets\min(x_i-l_i,\ g_i)$
  \EndIf
  \State $s\gets\max(s,|g_i|)$
\EndFor
\State \Return $s$
\end{algorithmic}
\caption{상자 제약의 사영 기울기 노름 (stardust \code{csrc/optimize/lbfgsb.cpp}의 \code{projgr})}
\label{alg:projgr}
\end{algorithm}

L-BFGS-B는 이 값이 \code{pgtol}(stardust의 \code{gtol}) 이하이면 멈춘다(\cref{alg:projgr}). 3행은 구현의 경계 처리이다. NaN과의 비교는 모두 거짓이므로 $\max(s,|g_i|)$에 NaN이 섞이면 \code{fmax}가 NaN을 버려 기울기가 NaN인데도 판정을 통과할 수 있다. 그래서 NaN을 만나면 그것을 그대로 돌려주어 문턱 비교가 거짓이 되게 한다. 경계에 정확히 붙은 성분이 $0$으로 사라지는 것은 상자 사영이 반올림 없는 비교라 $x_i=u_i$가 정확히 성립하기 때문이다.

\begin{figure}[htbp]
\centering
\begin{tikzpicture}[x=10mm,y=10mm,every node/.style={font=\small}]
\fill[black!10] (0,0) rectangle (4,2.6);
\draw[thick] (0,0) rectangle (4,2.6);
\node at (1.4,1.1) {$C$};
\fill[black!30] (4,2.6) -- (5.3,2.6) -- (5.3,3.8) -- (4,3.8) -- cycle;
\fill (4,2.6) circle (1.6pt);
\draw[->,thick] (4,2.6) -- (4.9,3.4) node[right] {$-\nabla f$};
\node at (4.95,2.85) {$N_C$};
\fill (2.2,2.6) circle (1.6pt);
\draw[->,thick] (2.2,2.6) -- (2.2,3.5) node[above] {$-\nabla f$};
\draw[->,dashed] (2.2,2.6) -- (2.95,3.2);
\node[anchor=west,font=\scriptsize] at (2.95,3.2) {최적 아님};
\node[anchor=north east] at (4,2.55) {꼭짓점};
\node[anchor=north] at (2.2,2.55) {변};
\end{tikzpicture}
\caption{상자의 법선 원뿔. 꼭짓점에서는 두 좌표의 음이 아닌 방향(진한 부분), 위쪽 변에서는 위쪽 방향 하나이다. $-\nabla f$가 법선 원뿔 안에 있으면 최적이고(\cref{eq:normal-cone-condition}) 변 위의 점에서 점선처럼 변을 따라가는 성분이 있으면 최적이 아니다. 사영 기울기는 바로 그 성분이다.}
\label{fig:normal-cone}
\end{figure}

\subsection{Functional Constraints and the KKT Theorem}

제약이 함수로 주어진 문제
\begin{equation}
\label{eq:nlp}
\min_x f(x)\quad\text{subject to}\quad c_i(x)\le0\ (i=1,\dots,m),\qquad Ax=b
\end{equation}
에서 법선 원뿔을 제약 함수의 기울기로 쓰면 KKT 조건이 된다. $c_i(x_\ast)=0$인 제약을 \term{활성}{active}이라 하고 그 첨자 집합을 $\mathcal A(x_\ast)$로 쓴다.

\begin{theorem}[KKT, 볼록 문제]
\label{thm:kkt-convex}
$f$와 $c_i$가 미분 가능한 볼록 함수이고 \term{Slater 조건}{Slater's condition}, 곧 $c_i(\bar x)<0$($i=1,\dots,m$)이고 $A\bar x=b$인 $\bar x$가 있다고 하자. 그러면 $x_\ast$가 \cref{eq:nlp}의 최소점일 필요충분조건은 다음을 만족하는 $\lambda\in\mathbb R^m$, $\nu$가 있는 것이다.
\[
\begin{gathered}
\nabla f(x_\ast)+\sum_i\lambda_i\nabla c_i(x_\ast)+A^{\mathsf T}\nu=0,\\
\lambda\ge0,\qquad c_i(x_\ast)\le0,\qquad Ax_\ast=b,\qquad\lambda_ic_i(x_\ast)=0.
\end{gathered}
\]
\end{theorem}

\begin{proof}
충분성: KKT 조건을 만족하는 $x_\ast$와 실현 가능한 $y$에 대해 볼록성에서 $f(y)\ge f(x_\ast)+\nabla f(x_\ast)^{\mathsf T}(y-x_\ast)$이고 $\lambda_i\ge0$, $c_i(y)\le0$, $\lambda_ic_i(x_\ast)=0$에서 $0\ge\lambda_ic_i(y)\ge\lambda_i\nabla c_i(x_\ast)^{\mathsf T}(y-x_\ast)$이며 $A(y-x_\ast)=0$이다. 셋을 합치면 $f(y)\ge f(x_\ast)+\nabla f(x_\ast)^{\mathsf T}(y-x_\ast)=f(x_\ast)-\sum_i\lambda_i\nabla c_i(x_\ast)^{\mathsf T}(y-x_\ast)-\nu^{\mathsf T}A(y-x_\ast)\ge f(x_\ast)$이다.

필요성은 다음 절의 강쌍대성(\cref{thm:strong-duality})에서 나온다. 강쌍대성이 주는 쌍대 최적해 $(\lambda,\nu)$에 대해 $f(x_\ast)=\inf_xL(x,\lambda,\nu)\le L(x_\ast,\lambda,\nu)\le f(x_\ast)$이므로 둘째 부등식의 등호에서 $\lambda_ic_i(x_\ast)=0$이고, 첫째 부등식의 등호에서 $x_\ast$는 볼록 함수 $L(\cdot,\lambda,\nu)$의 최소점이라 그 기울기가 $0$이다.
\end{proof}

제약이 볼록이 아니면 Slater 조건 대신 활성 제약의 기울기와 등식 제약의 기울기가 일차독립이라는 LICQ 같은 \term{제약 자격}{constraint qualification}이 필요하다. 그러면 KKT는 국소 최소점의 필요조건이 된다. 증명은 먼저 활성 제약의 선형화가 만드는 원뿔이 접원뿔과 같음을 보이고(여기서 LICQ와 음함수 정리를 쓴다) \cref{eq:normal-cone-condition}의 꼴을 그 원뿔에 쓴 뒤 \cref{lem:farkas}로 승수의 존재를 끌어낸다(Nocedal과 Wright\cite{nocedal2006} 12.3절). 조건의 이름은 세 사람에게서 왔다. Karush\cite{karush1939}가 1939년 석사 논문에서 처음 썼고 Kuhn과 Tucker\cite{kuhn1951}가 1951년에 다시 발견해 널리 알렸다. 제약 자격이 왜 필요한지는 간단한 예가 보여 준다. $\min x$ subject to $x^2\le0$의 유일한 실현 가능 점 $0$은 최소점이지만 $1+\lambda\cdot2x=1\ne0$이어서 승수가 없다. Slater 조건이 깨져 있기 때문이다. 제약이 있으면 이계 조건은 임계 원뿔 위에서 따진다. 임계 원뿔은 KKT 점에서 활성 제약을 깨뜨리지 않으면서 함수값을 일계까지 바꾸지 않는 방향들의 원뿔이고 그 위에서 $d^{\mathsf T}\nabla^2_{xx}L\,d\ge0$이 필요하고 $>0$이면 엄격한 국소 최소점이다. \cref{thm:optimality}를 임계 원뿔을 따라가는 곡선에 쓰면 증명된다(Nocedal과 Wright\cite{nocedal2006} 정리 12.5, 12.6). 상자 제약과 비음 최소제곱에서는 이 조건이 엄격한 상보성(묶인 변수의 승수가 양수)과 자유 변수 열들의 독립으로 나타나며 둘이 함께 해의 유일성을 보증한다.

상자 제약 $l\le x\le u$는 $c(x)=l-x$, $x-u$로 쓸 수 있고 KKT 조건은 $\nabla f(x)=\lambda^l-\lambda^u$, $\lambda^l,\lambda^u\ge0$, 상보성이 된다. 이것을 성분마다 풀면 바로 앞 절의 ``사영 기울기가 $0$''이다. \code{lsq\_linear}이 결과로 주는 \code{optimality}는 이 KKT 잔차를 방법마다 다른 척도로 잰 값이다. \code{'bvls'}는 묶인 변수에서 $g_i$에 경계 쪽 부호를 곱한 것과 자유 변수의 $|g_i|$ 가운데 최댓값이고 \code{'trf'}는 Coleman--Li의 척도 $v_i$를 곱한 $\|v\cdot g\|_\infty$이다. 척도가 다르므로 두 방법의 \code{optimality}는 같은 문제에서도 다른 값이 나온다(\cref{tab:box-kkt}). \code{'bvls'}의 값은 음수가 될 수도 있다. $A=I$, $b=(-1,-2,5)$, 상자 $[0,1]^3$에서 해는 $(0,0,1)$이고 세 변수가 모두 묶여 엄격한 상보성의 여유가 그대로 나오므로 \code{optimality}가 $-1$이었다. 같은 문제를 \code{'trf'}로 풀면 해가 $(4.1\times10^{-11},\,5\times10^{-324},\,1-2^{-53})$으로 경계에 붙지 않는다. 신뢰 영역 반사법은 반복점을 늘 엄격한 내부에 둔다. 그래서 \code{'trf'}의 \code{active\_mask}는 경계까지의 상대 거리가 \code{tol}$\cdot\max(1,|\text{경계}|)$ 이하인 변수로 정한다(\code{figures/08/bvls\_neg.py}).

\section{Lagrange Duality}
\label{sec:duality}

\subsection{Weak and Strong Duality}

\cref{eq:nlp}의 \term{Lagrange 함수}{Lagrangian}
\[
L(x,\lambda,\nu)=f(x)+\sum_i\lambda_ic_i(x)+\nu^{\mathsf T}(Ax-b)
\]
에서 $x$에 대한 하한 $g(\lambda,\nu)=\inf_xL(x,\lambda,\nu)$를 \term{쌍대 함수}{dual function}라 한다. 일차 함수들의 하한이므로 원래 문제가 볼록이 아니어도 $g$는 오목이다.

\begin{proposition}[약쌍대성]
\label{prop:weak-duality}
$\lambda\ge0$이고 $x$가 실현 가능하면 $g(\lambda,\nu)\le f(x)$이다. 따라서 \term{쌍대 간격}{duality gap} $f(x)-g(\lambda,\nu)$는 음이 아니고 $f(x)-f_\ast$의 상한이다.
\end{proposition}

\begin{proof}
실현 가능한 $x$에서 $\lambda_ic_i(x)\le0$, $Ax-b=0$이므로 $g(\lambda,\nu)\le L(x,\lambda,\nu)\le f(x)$이다. 이것이 모든 실현 가능한 $x$에서 성립하므로 $g(\lambda,\nu)\le f_\ast$이다.
\end{proof}

계산에서는 이 단순한 부등식이 요긴하다. 반복 중의 $x$와 아무 쌍대 실현 가능 점 $(\lambda,\nu)$를 들고 있으면 간격이 ``$x$가 최적값에서 얼마나 떨어져 있는지''를 위에서 막는 증명서가 된다. 기울기 판정과 달리 이 수는 목적 함수와 같은 단위를 갖는다.

\begin{theorem}[강쌍대성]
\label{thm:strong-duality}
\cref{thm:kkt-convex}의 가정(볼록, Slater 조건) 아래에서 최적값 $f_\ast$가 유한하면 $g(\lambda,\nu)=f_\ast$인 $\lambda\ge0$, $\nu$가 있다.
\end{theorem}

\begin{proof}
줄거리를 적는다(Boyd와 Vandenberghe\cite{boyd2004convex} 5.3.2절). 집합
\[
\mathcal B=\{(u,w,t):\exists x,\ c_i(x)\le u_i,\ Ax-b=w,\ f(x)\le t\}
\]
은 볼록이고 $(0,0,f_\ast)$는 그 내부에 있지 않다($t<f_\ast$로 갈 수 없다). \cref{thm:separation}의 지지 초평면에서 $(\lambda,\nu,\mu)\ne0$이 있어 $\mathcal B$의 모든 점에서 $\lambda^{\mathsf T}u+\nu^{\mathsf T}w+\mu t\ge\mu f_\ast$이다. $u$와 $t$를 키울 수 있으므로 $\lambda\ge0$, $\mu\ge0$이다. $\mu=0$이라면 Slater 점 $\bar x$에서 $\lambda^{\mathsf T}c(\bar x)\ge0$이고 $c(\bar x)<0$이므로 $\lambda=0$, 그러면 $\nu^{\mathsf T}(Ax-b)\ge0$이 모든 $x$에서 성립해 $A^{\mathsf T}\nu=0$이다. 여기서 $A$의 행이 일차독립이라고 두어도 일반성을 잃지 않는다. Slater 점이 $Ax=b$를 만족하므로 종속인 행은 남은 행들의 결합이고 그 행을 빼도 실현 가능 집합과 $f_\ast$가 그대로이며 뺀 행의 승수는 $0$으로 두면 되기 때문이다. 그러면 $\nu=0$이 되어 $(\lambda,\nu,\mu)\ne0$에 어긋난다. 따라서 $\mu>0$이고 $\mu=1$로 나누면 모든 $x$에서 $L(x,\lambda,\nu)\ge f_\ast$, 곧 $g(\lambda,\nu)\ge f_\ast$이다. 약쌍대성과 합쳐 등호이다.
\end{proof}

약쌍대성은 간격이 작으면 최적에 가깝다는 것을 보증한다. 최적에서 간격을 $0$으로 만드는 쌍대 점이 있다는 것은 강쌍대성이 보증한다. 증명에 나온 초평면에서는 승수의 뜻도 읽을 수 있다.

\begin{corollary}[승수의 감도]
\label{cor:sensitivity}
\cref{thm:strong-duality}의 가정 아래에서 제약을 늦춘 문제의 최적값을 $p^\ast(u,w)=\inf\{f(x):c(x)\le u,\ Ax-b=w\}$로 쓰면 쌍대 최적해 $(\lambda^\ast,\nu^\ast)$에 대해 모든 $(u,w)$에서
\[
p^\ast(u,w)\ge p^\ast(0,0)-\lambda^{\ast\mathsf T}u-\nu^{\ast\mathsf T}w
\]
이다. 따라서 $p^\ast$가 $(0,0)$에서 미분 가능하면 $\lambda^\ast=-\nabla_up^\ast(0,0)$이다.
\end{corollary}

\begin{proof}
$c(x)\le u$, $Ax-b=w$인 $x$에서 $p^\ast(0,0)=g(\lambda^\ast,\nu^\ast)\le f(x)+\lambda^{\ast\mathsf T}c(x)+\nu^{\ast\mathsf T}(Ax-b)\le f(x)+\lambda^{\ast\mathsf T}u+\nu^{\ast\mathsf T}w$이고 이런 $x$에 대해 하한을 취한다.
\end{proof}

승수 $\lambda_i^\ast$는 제약 $i$를 조금 늦출 때 최적값이 줄어드는 빠르기이며 그 제약의 그림자 가격이라고도 부른다. \cref{sec:nnls}의 비음 최소제곱에서 $0$에 묶인 계수의 $-w_j$가 그 하한 제약의 승수이고 라소에서 최적값을 $\lambda_1$로 미분하면 $\|b_\ast\|_1$이다(포락선 정리). 이것이 쌍대 간격을 멈춤 판정으로 쓰는 근거이다.

\subsection{The Lasso Dual and Its Gap}

라소 \cref{eq:lasso}($\lambda_2=0$)의 쌍대는 잔차를 변수로 따로 두면 깔끔하게 나온다. $\min(1/2)\|r\|_\Omega^2+\lambda_1\|b\|_1$ subject to $r=v-Zb$에서 승수를 $\Omega\theta$로 두면
\[
L=(1/2)\|r\|_\Omega^2+\lambda_1\|b\|_1+\theta^{\mathsf T}\Omega(v-Zb-r)
\]
이다. $r$에 대한 하한은 $r=\theta$에서 $-(1/2)\|\theta\|_\Omega^2$이고 $b$에 대한 하한은 켤레 함수의 계산에서 $\|Z^{\mathsf T}\Omega\theta\|_\infty\le\lambda_1$이면 $0$, 아니면 $-\infty$이다. 따라서 쌍대 문제는
\begin{equation}
\label{eq:lasso-dual}
\begin{aligned}
&\max_\theta\ D(\theta)=\theta^{\mathsf T}\Omega v-\frac12\|\theta\|_\Omega^2=\frac12\|v\|_\Omega^2-\frac12\|v-\theta\|_\Omega^2\\
&\text{subject to}\quad\|Z^{\mathsf T}\Omega\theta\|_\infty\le\lambda_1
\end{aligned}
\end{equation}
이다. 최적에서는 $\theta_\ast=r_\ast$이고 쌍대 제약은 \cref{eq:lasso-kkt}의 $|z_j^{\mathsf T}\Omega r_\ast|\le\lambda_1$과 같다. 반복 중의 잔차 $r$은 대개 쌍대 제약을 조금 어기므로 척도를 줄여 실현 가능하게 만든다. $s=\min(1,\lambda_1/\|Z^{\mathsf T}\Omega r\|_\infty)$로 두면 $\theta=sr$은 실현 가능하고 \cref{prop:weak-duality}의 간격 $P(b)-D(sr)$가 $P(b)-P_\ast$의 상한이다.

이 간격을 정의 그대로 계산하면 문제가 생긴다. $D(sr)=(1/2)\|v\|_\Omega^2-(1/2)\|v-sr\|_\Omega^2$는 $(1/2)\|v\|_\Omega^2$ 크기의 두 항의 차이고 해에 가까워지면 이렇게 얻은 $D$를 다시 $P$에서 빼서 작은 간격을 만든다. stardust는 식을 미리 정리해 그런 항이 서로 지워지지 않게 쓴다. $R^2=\|r\|_\Omega^2$, $D_f=r^{\mathsf T}\Omega(v-r)$(맞춘 값과 잔차의 곱)으로 두면 $\|v-sr\|_\Omega^2=\|v\|_\Omega^2-2s\,r^{\mathsf T}\Omega v+s^2R^2$이고 $r^{\mathsf T}\Omega v=R^2+D_f$이므로
\begin{equation}
\label{eq:gap-stable}
P(b)-D(sr)=\frac12R^2(1-s)^2+\lambda_1\|b\|_1-sD_f
\end{equation}
이다. 남은 항들은 모두 잔차나 계수의 크기이고 $\|v\|_\Omega^2$이 사라졌다. 엘라스틱넷($\lambda_2>0$)은 $Z$ 아래에 $\sqrt{\lambda_2}I$를 쌓고 $v$ 아래에 $0$을 붙인 라소이다(붙인 부분의 무게는 $1$). 늘린 잔차 $(r,-\sqrt{\lambda_2}b)$에 같은 식을 쓰면 $R^2$가 $R^2+\lambda_2\|b\|^2$로, $D_f$가 $D_f-\lambda_2\|b\|^2$로 바뀌어 $(1/2)R^2(1-s)^2+(1/2)\lambda_2\|b\|^2(1+s^2)+\lambda_1\|b\|_1-sD_f$를 얻고 쌍대 제약의 내적은 $Z^{\mathsf T}\Omega r-\lambda_2b$가 된다.

\begin{algorithm}[htbp]
\begin{algorithmic}[1]
\Require 계수 $b$, 잔차 $r$(또는 Gram 행렬에서 다시 만든 $R^2$과 $c-Gb$), $\lambda_1$, $\lambda_2$
\Ensure 쌍대 간격 $P(b)-D(sr)$
\State $R^2\gets\|r\|_\Omega^2$; $\xi\gets Z^{\mathsf T}\Omega r$; $D_f\gets b^{\mathsf T}\xi$ \Comment{절편이 있으면 그 항을 더한다}
\State $\xi\gets\xi-\lambda_2b$; $\delta\gets\max_j|\xi_j|$ (상수 열은 뺀다)
\State $s\gets\lambda_1/\delta$ ($\delta>\lambda_1$), 아니면 $1$ \Comment{$sr$을 쌍대 실현 가능하게}
\State \Return $(1/2)R^2(1-s)^2+(1/2)\lambda_2\|b\|^2(1+s^2)+\lambda_1\|b\|_1-sD_f$ \Comment{\cref{eq:gap-stable}}
\end{algorithmic}
\caption{라소·엘라스틱넷의 소거 없는 쌍대 간격 (stardust \code{penalized.cpp}의 \code{gap}, \code{gap\_from})}
\label{alg:gap}
\end{algorithm}

stardust는 간격이 \code{tol}$\cdot\|v\|_\Omega^2$ 이하이면 수렴으로 본다(scikit-learn과 같은 기준). \cref{fig:lasso-gap}은 두 식의 차이를 잰 것이다. stardust의 라소 해 $b_\ast$를 $b_\ast+\epsilon d$로 흔들어 간격을 정확한 유리수로 계산하고 binary64의 두 식과 견주었다. 반응이 큰 신호를 담아 $\|v\|_\Omega^2$이 목적 함수 값의 약 $7700$배인 예이다. 정의대로 쓴 식은 상대 간격 $4\times10^{-13}$ 근처에서 바닥에 닿는다. 그 아래로는 간격을 읽지 못하고 부호까지 틀린다. \cref{eq:gap-stable}은 정확한 간격을 $4\times10^{-15}$까지 따라간다. 해 자체에서 stardust가 보고한 간격 $3.75\times10^{-12}$은 정확한 간격 $3.72\times10^{-12}$과 맞는다. 이 바닥 아래의 간격은 해를 binary64로 반올림한 데서 오므로 더 줄일 수 없다. \cref{eq:gap-stable}에도 소거가 아주 사라지지는 않는다. 해 근처에서 $\lambda_1\|b\|_1-sD_f$는 크기가 $P$ 이하인 두 항의 차이므로 바닥이 $u\|v\|_\Omega^2$에서 $uP$로 내려간 것이고 두 바닥의 비가 바로 $\|v\|_\Omega^2/P\approx7700$이다. 그림의 점들은 잔차 $r$로 직접 계산한 값이다. 열이 $1024$개 이하인 이 문제에서 stardust는 Gram 경로를 타서 $R^2=v^{\mathsf T}\Omega v-2b^{\mathsf T}c+b^{\mathsf T}Gb$와 $\xi=c-Gb$($c=Z^{\mathsf T}\Omega v$)를 쓰는데 첫 식에 소거가 있어 $R^2$의 상대오차가 $10^{-9}$ 안팎이다. 다행히 $R^2$에는 $(1-s)^2$이 곱해지므로 간격에 미치는 영향은 작고 같은 점들을 그 경로로 다시 재면 상대 간격 $10^{-13}$까지는 정확한 값을 따라가다가 $1.0\times10^{-15}$에서 바닥에 닿았다(정확한 값 $4.1\times10^{-15}$; \code{figures/08/gap\_gram.py}). 간격은 가장 큰 계수 변화가 \code{tol}$\cdot$(가장 큰 계수) 이하인 전체 쓸기 뒤에만 계산한다.

\begin{figure}[htbp]
\centering
\begin{tikzpicture}
\begin{loglogaxis}[sdplot,xlabel={흔든 크기 $\epsilon$},ylabel={상대 쌍대 간격},xmin=1e-16,xmax=1,ymin=1e-16,ymax=1e-1,x dir=reverse]
\addplot[black,thick] table[x=eps,y=exact]{figures/08/lasso_gap.dat};
\addlegendentry{정확한 간격}
\addplot[black,only marks,mark=o,mark size=1.6pt] table[x=eps,y=naive]{figures/08/lasso_gap.dat};
\addlegendentry{$P-D$ 그대로}
\addplot[black,only marks,mark=*,mark size=1pt] table[x=eps,y=stable]{figures/08/lasso_gap.dat};
\addlegendentry{\cref{eq:gap-stable}}
\end{loglogaxis}
\end{tikzpicture}
\caption{라소의 쌍대 간격을 목적 함수 값으로 나눈 것(따로 빌드한 stardust의 해 $b_\ast$ 근처, $n=40$, $p=12$, $\|v\|_\Omega^2/P\approx7700$; 정확한 값은 유리수 계산, \code{figures/08/lasso\_gap.py}). 빈 원은 $|P-D|$를 정의대로 계산한 값이다.}
\label{fig:lasso-gap}
\end{figure}

간격만으로 판정하면 좌표 하강은 활성 집합의 조건이 나쁠 때 선형으로 느리게 수렴한다. stardust는 이 약점을 KKT 조건으로 메운다. 쓸기 두 번이 계수의 부호를 바꾸지 않으면 활성 집합 $\mathcal A$와 부호 $\sigma$가 자리 잡았다고 보고 \cref{eq:lasso-kkt}를 그 부호로 고정해 선형계
\[
(Z_{\mathcal A}^{\mathsf T}\Omega Z_{\mathcal A}+\lambda_2I)\,b_{\mathcal A}=Z_{\mathcal A}^{\mathsf T}\Omega v-\lambda_1\sigma
\]
를 Cholesky로 푼다(\cref{alg:polish}). 이 해가 부호를 지키고 비활성 열마다 $|z_j^{\mathsf T}\Omega r|\le\lambda_1$이면 \cref{eq:lasso-kkt}를 모두 만족하므로 정확한 최소점이다(볼록 문제이므로 \cref{prop:subgradient-optimality}). 이 ``다듬기''의 결과는 $\mathcal A$와 $\sigma$에만 의존하므로 경로로 풀든 한 번에 풀든 같은 비트가 나온다.

\begin{algorithm}[htbp]
\begin{algorithmic}[1]
\Require 계수 $b$, Gram $G=Z^{\mathsf T}\Omega Z$, $c=Z^{\mathsf T}\Omega v$, $\lambda_1$, $\lambda_2$
\Ensure 다듬은 $b$ 또는 실패
\State $\mathcal A\gets\{j:b_j\ne0\}$; $\sigma\gets\operatorname{sign}(b_{\mathcal A})$ \Comment{복제 열은 첫 열 하나만, 합의 부호로}
\State $H\gets G_{\mathcal A\mathcal A}+\lambda_2I$; $L\gets\operatorname{chol}(H)$ (실패하면 실패)
\State $b_{\mathcal A}\gets H^{-1}(c_{\mathcal A}-\lambda_1\sigma)$
\If{$1/\operatorname{rcond}_1(L)^2>64|\mathcal A|$} \Comment{조건수 추정이 크면}
  \For{$k=1,2,3$} \Comment{자료에서 double-double로 만든 잔차}
    \State $\rho\gets Z_{\mathcal A}^{\mathsf T}\Omega(v-Z_{\mathcal A}b_{\mathcal A})-\lambda_2b_{\mathcal A}-\lambda_1\sigma$; $\delta\gets H^{-1}\rho$; $b_{\mathcal A}\gets b_{\mathcal A}+\delta$
    \State \textbf{if} $\|\delta\|_\infty\le4u\|b_{\mathcal A}\|_\infty$ \textbf{or} $\|\delta\|_\infty$가 반으로 줄지 않았다 \textbf{then} \textbf{break}
  \EndFor
\EndIf
\State \textbf{if} 어떤 $q\in\mathcal A$에서 $\sigma_qb_q\le0$ \textbf{then} 실패 \Comment{부호가 지켜져야 한다}
\For{$j\notin\mathcal A$, $d_j>0$}
  \State \textbf{if} $|c_j-G_{j\mathcal A}b_{\mathcal A}|>\lambda_1(1+10^{-9})+64u\sqrt{d_j\,v^{\mathsf T}\Omega v}$ \textbf{then} 실패
\EndFor
\State \Return $b$ ($\mathcal A$ 밖은 $0$)
\end{algorithmic}
\caption{라소·엘라스틱넷의 KKT 다듬기 (stardust \code{penalized.cpp}의 \code{polish})}
\label{alg:polish}
\end{algorithm}

비활성 열의 판정에 붙은 여유는 \cref{alg:cd}의 $\epsilon_j$와 같은 크기에 상대 여유 $10^{-9}$을 더한 것이라 쓸기가 $0$에 묶어 둔 열은 다듬기도 받아들인다. 그 여유가 활성 집합을 잘못 고르게 할 수 있는지 따로 빌드한 stardust로 시험했다. 새 열의 KKT 위반을 $10^{-6}$부터 $10^{-12}$까지 줄여 가며 만들었을 때 모두 그 열을 들였다. 다듬기의 비용은 $|\mathcal A|^3/3+p|\mathcal A|$이고 반복 개선을 붙이면 $O(m|\mathcal A|)$가 더해진다. 다듬기는 바깥 고리(\cref{alg:cd-outer})에서 하강 중에 부호 패턴이 두 번 연속 같을 때(패턴마다 한 번)와 간격 판정을 통과한 뒤에 시도한다. 뒤의 경우 다듬기가 실패하면 \code{tol}을 $10^{-6}$배로 줄여 좌표 하강을 두 번까지 다시 한다.

\begin{algorithm}[htbp]
\begin{algorithmic}[1]
\Require $b$, $\lambda_1$, $\lambda_2$, \code{tol}, 반복 상한
\Ensure 해와 수렴 여부
\State $\tau\gets\mathrm{tol}\cdot\|v\|_\Omega^2$; \textbf{if} \Call{Gap}{$b$}$\le\tau$ \textbf{then} \Return 수렴
\Loop
  \State 모든 열을 쓸고(\cref{alg:cd}) $\Delta_{\max}$, 가장 큰 계수 $w_{\max}$를 얻는다
  \State \textbf{if} \Call{Settle}{} \textbf{then} \Return 수렴(다듬기)
  \If{$w_{\max}=0$ \textbf{or} $\Delta_{\max}\le\mathrm{tol}\cdot w_{\max}$ \textbf{or} 상한}
    \State \textbf{if} \Call{Gap}{$b$}$\le\tau$ \textbf{then} \Return 수렴 \Comment{간격은 이때만, \cref{alg:gap}}
    \State \textbf{if} 상한 \textbf{then} \Return 미수렴
  \EndIf
  \Repeat \Comment{0이 아닌 계수만 도는 짧은 쓸기}
    \State $\mathcal A\gets\{j:b_j\ne0\}$에서 쓸고 \Call{Settle}{}이 참이면 \Return 수렴(다듬기)
  \Until{$\Delta_{\max}\le\mathrm{tol}\cdot w_{\max}$ \textbf{or} 상한}
\EndLoop
\Function{Settle}{} \Comment{부호 패턴이 자리 잡았으면 다듬기를 한 번 시도}
  \State $\sigma\gets\operatorname{sign}(b)$; \textbf{if} $\sigma=\sigma_{\mathrm{prev}}\ne\sigma_{\mathrm{tried}}$: $\sigma_{\mathrm{tried}}\gets\sigma$, \Call{Polish}{}이 성공하고 \Call{Gap}{}$\le\tau$이면 참
  \State $\sigma_{\mathrm{prev}}\gets\sigma$; 거짓
\EndFunction
\end{algorithmic}
\caption{좌표 하강의 바깥 고리 (stardust \code{penalized.cpp}의 \code{cd\_solve})}
\label{alg:cd-outer}
\end{algorithm}

쓸기 한 번의 비용은 Gram 경로에서 열마다 $O(1)$에 계수가 바뀐 열의 $O(p)$ 갱신이고 잔차 경로에서 $O(m)$이다. 간격은 Gram 경로에서 $O(p^2)$, 잔차 경로에서 $O(mp)$이다. 행에 걸친 합(평균, Gram, 잔차)은 크기로만 정한 $4096$행 조각으로 나누어 조각의 부분합을 차례로 더하므로 스레드 수와 상관없이 같은 비트가 나온다. 경로와 교차검증을 포함한 벌점 회귀 전체는 126장에서 다룬다. 분위 회귀(128장)에서도 쌍대가 멈춤 판정을 맡는다. 분위 회귀는 선형 계획의 쌍대 문제를 내부점법으로 풀며 상대 쌍대 간격이 \code{tol} 아래로 내려가면 멈추고 이어지는 심플렉스 걸음은 승수 $\lambda_j$가 모두 상자 $[m_j(q-1),m_jq]$ 안에 드는 쌍대 증명서가 설 때까지 돈다.

\section{Active Sets: Nonnegative Least Squares}
\label{sec:nnls}

\subsection{Optimality Conditions}

KKT 조건에서 알고리즘을 직접 설계하는 예로 \term{비음 최소제곱}{nonnegative least squares}
\begin{equation}
\label{eq:nnls}
\min_x\ (1/2)\|Ax-b\|^2\quad\text{subject to}\quad x\ge0
\end{equation}
을 보자. 스펙트럼의 혼합 비율, 농도, 분산 성분처럼 음수일 수 없는 계수를 맞출 때 쓴다. 볼록 문제이고 제약 $-x_j\le0$은 Slater 조건을 만족하므로 \cref{thm:kkt-convex}에서 $x$가 최적일 필요충분조건은 $w=A^{\mathsf T}(b-Ax)$(음의 기울기)에 대해
\begin{equation}
\label{eq:nnls-kkt}
x\ge0,\qquad w\le0\ \text{on}\ \{j:x_j=0\},\qquad w_j=0\ \text{on}\ \{j:x_j>0\}
\end{equation}
이다. 양수인 계수의 집합 $P$를 알면 문제는 $P$의 열만 쓰는 보통의 최소제곱이 된다. 그래서 남은 일은 $P$를 찾는 것뿐이며 그 집합을 어떻게 고칠지는 \cref{eq:nnls-kkt}가 알려 준다. $w_j>0$인 $j\notin P$는 늘리면 목적 함수가 주는 방향이므로 $P$에 들일 후보이다.

\begin{lemma}
\label{lem:nnls-enter}
$x$가 열 집합 $P$ 위의 최소제곱 해이고 $r=b-Ax$에 대해 $t\notin P$에서 $w_t=a_t^{\mathsf T}r>0$이면 $P\cup\{t\}$ 위의 최소제곱 해 $z$의 $t$번째 성분은 양수이다.
\end{lemma}

\begin{proof}
$\Pi$를 $P$의 열들이 생성한 공간으로의 직교 사영이라 하자. $x$가 그 공간의 최소제곱 해이므로 $r\perp\operatorname{span}A_P$, 곧 $\Pi r=0$이다. $\tilde a_t=(I-\Pi)a_t$로 두면 $\tilde a_t^{\mathsf T}r=a_t^{\mathsf T}r=w_t>0$이므로 $\tilde a_t\ne0$이다. 열 하나를 더한 최소제곱의 정규방정식은 $A_P^{\mathsf T}(A_Pz_P+a_tz_t-b)=0$과 $a_t^{\mathsf T}(A_Pz_P+a_tz_t-b)=0$이다. 첫 식에서 $A_Pz_P=\Pi(b-a_tz_t)$이고 이것을 둘째 식에 넣으면 $a_t^{\mathsf T}(I-\Pi)(a_tz_t-b)=0$, 곧 $z_t\tilde a_t^{\mathsf T}\tilde a_t=\tilde a_t^{\mathsf T}b=\tilde a_t^{\mathsf T}r$이다(Frisch--Waugh). 따라서 $z_t=w_t/\|\tilde a_t\|^2>0$이다.
\end{proof}

\subsection{The Lawson--Hanson Algorithm}

Lawson과 Hanson\cite{lawson1995}의 활성 집합법(\cref{alg:nnls})은 바깥 고리에서 열을 하나 들이고 안쪽 고리에서 실현 가능성을 되찾는다. $P$ 위의 최소제곱 해 $z$가 모두 양수이면 그대로 받아들인다. 음이 되거나 $0$이 된 성분이 있으면 $x$에서 $z$로 가는 선분을 따라 처음으로 경계에 닿는 곳까지만 간다. 그 걸음 $\alpha=\min\{x_j/(x_j-z_j):j\in P,\ z_j\le0\}$은 $(0,1]$에 있다. $P$의 성분은 모두 $x_j>0$이므로 각 비는 양수이고 $z_j\le0$이므로 $1$ 이하이며 $1$은 $z_j=0$인 성분이 최소를 이룰 때만 나온다. 그 점에서 $0$이 된 계수를 $P$에서 빼고 다시 푼다.

\begin{algorithm}[htbp]
\begin{algorithmic}[1]
\Require $A\in\mathbb R^{m\times n}$, $b$, 반복 상한
\Ensure \cref{eq:nnls}의 해 $x$
\State $P\gets\emptyset$; $x\gets0$
\Loop
  \State $w\gets A^{\mathsf T}(b-Ax)$ \Comment{반사한 행렬의 아래 부분으로 $O(m-|P|)$}
  \State \textbf{if} $|P|=\min(m,n)$ \textbf{or} $P$ 밖의 모든 $j$에서 $w_j\le0$ \textbf{then} \Return $x$ \Comment{\cref{eq:nnls-kkt}}
  \State $t\gets\arg\max_{j\notin P}w_j$; $P\gets P\cup\{t\}$ \Comment{\cref{lem:nnls-enter}: $z_t>0$}
  \Loop
    \State $z\gets\arg\min\{\|A_Pz_P-b\|:z_j=0\ (j\notin P)\}$ \Comment{QR을 열 하나씩 고친다}
    \State \textbf{if} 모든 $j\in P$에서 $z_j>0$ \textbf{then} $x\gets z$; \textbf{break}
    \State $\alpha\gets\min\{x_j/(x_j-z_j):j\in P,\ z_j\le0\}$; $x\gets x+\alpha(z-x)$
    \State $x_j\le0$인 $j$를 $P$에서 빼고 $x_j\gets0$; $P=\emptyset$이면 \textbf{break}
  \EndLoop
\EndLoop
\end{algorithmic}
\caption{Lawson--Hanson의 비음 최소제곱 (stardust \code{csrc/optimize/linear.cpp}의 \code{nnls\_kernel})}
\label{alg:nnls}
\end{algorithm}

\begin{theorem}
\label{thm:nnls-finite}
정확한 산술에서 \cref{alg:nnls}는 유한 번 만에 \cref{eq:nnls-kkt}를 만족하는 $x$에서 멈춘다.
\end{theorem}

\begin{proof}
줄거리를 적는다(Lawson과 Hanson\cite{lawson1995} 23장). 먼저 $A_P$의 열은 언제나 일차독립이다. 들이는 열은 $\tilde a_t\ne0$(\cref{lem:nnls-enter})이므로 독립을 깨지 않고 빼는 것은 독립을 지킨다. 그래서 $P$ 위의 최소제곱 해는 하나로 정해진다. 안쪽 고리는 매번 $P$에서 열을 적어도 하나 빼므로 많아야 $|P|$번 돌고 끝나면 $x$는 지금의 $P$ 위의 최소제곱 해이며 모든 성분이 양수이다. 바깥 고리가 열 $t$를 들일 때 \cref{lem:nnls-enter}에서 새 해의 $t$ 성분이 양수이므로 안쪽 고리가 끝난 뒤의 잔차는 들이기 전보다 엄격히 작다($\phi(x)=(1/2)\|Ax-b\|^2$로 두자. 들이기 직전의 $x$는 옛 $P$ 위의 해이고 $w_t\ne0$이므로 새 $P$ 위의 해 $z$에서 $\phi(z)<\phi(x)$이다. 첫 안쪽 걸음의 $\alpha$는 양수이므로 볼록성에서 $\phi(x+\alpha(z-x))\le(1-\alpha)\phi(x)+\alpha\phi(z)<\phi(x)$이고 그 뒤의 걸음들도 $\phi$를 늘리지 않는다). 바깥 고리가 끝날 때의 $x$는 그때의 $P$로 정해지는 최소제곱 해이므로 같은 $P$가 두 번 나오면 잔차가 같아야 하는데 잔차가 엄격히 줄었으므로 그럴 수 없다. 부분집합은 $2^n$개뿐이므로 고리는 끝나고, 끝나는 조건이 바로 \cref{eq:nnls-kkt}이다.
\end{proof}

\subsection{Implementation}

stardust의 \code{optimize.nnls}는 scipy 1.18.1의 C 판(Lawson--Hanson을 LAPACK으로 다시 쓴 것)을 기계어에서 읽어 옮기되 짧은 내적만 직접 짠 고리로 바꾸었다. 자료 배치와 수치 처리는 다음과 같다.
\begin{enumerate}
\item $P$의 열들은 열 순서를 바꾼 $A$의 앞쪽에 모이고 그 부분은 Householder 반사로 위 삼각이 된 상태를 유지한다. 열을 들일 때는 그 열의 $P$ 아래 부분에 반사 하나(LAPACK \code{dlarfgp}: 대각 $\beta$가 음이 아닌 꼴)를 만들어 나머지 열과 $b$에 적용하고 7행의 $z$는 삼각계의 뒤 대입 한 번이다. 열을 뺄 때는 그 뒤의 열들을 한 칸씩 당기고 생긴 아래 대각 원소를 Givens 회전(\code{dlartgp})으로 없앤다. 이렇게 하면 바깥 고리 한 번의 비용이 처음부터 QR을 다시 하는 $O(mn^2)$ 대신 $O(mn)$이다.
\item 들이기 판정이 둘 더 있다. 새 열의 $P$ 밖 부분의 크기 $|\beta|$가 그 열의 $P$ 부분 노름의 ulp의 $100$배보다 작으면 그 열은 이미 $P$의 열들로 거의 표현되므로(\cref{lem:nnls-enter}의 $\tilde a_t\approx0$) 들이지 않는다. 또 반사를 적용한 $b$에서 새 계수의 부호 $z_t>0$을 먼저 확인하고 그렇지 않으면 $w_t=0$으로 두고 다음 후보로 넘어간다. 정확한 산술이라면 \cref{lem:nnls-enter}가 보증할 부호를 반올림이 뒤집는 경우를 이 확인이 막는다.
\item 9행의 $\alpha$는 $2$에서 시작한다. 가능한 걸음은 모두 $1$ 이하이므로 $\alpha$가 그대로 $2$이면 ``$z$가 실현 가능하다''는 표지이고 그 경우는 8행과 같다. 갱신은 \code{fma}로 한다.
\item 3행의 쌍대 값 $w$는 $A^{\mathsf T}(b-Ax)$로 계산하지 않는다. 반사를 적용한 $A$와 $b$에서 위쪽 $|P|$행의 잔차는 $P$ 위의 최소제곱 해에서 정확히 $0$이므로 $w_j$는 열 $j$와 $b$의 $|P|$행 아래 부분의 내적이고 비용이 열마다 $O(m-|P|)$이다.
\item 내적과 열 반사는 길이가 $256$ 미만이면 누산기 네 개의 고정된 순서로, 그 이상이면 BLAS로 한다. 짧은 길이에서는 Accelerate를 부르는 비용이 계산 자체보다 커서 이렇게 바꾸었고 $(100,10)$에서 약 $1.3$배, $(10^3,100)$에서 $2$배 남짓 빨라졌다(설계 문서의 측정). 그 대가로 scipy와 비트까지 같지는 않다. 해의 $0$인 자리는 같고 나머지는 부분 문제의 조건수에 비례하는 끝자리 차이가 난다. 합의 순서가 길이로만 정해지므로 어느 열을 들일지는 실행마다 같다.
\item 바깥 고리는 $|P|=\min(m,n)$이 되거나 후보 열이 모두 들이기 판정에서 떨어져 더 들일 열이 없으면 끝난다. 10행에서 빼는 열은 $\alpha$로 $0$이 된 열 하나만이 아니라 반올림으로 $x_j\le0$이 된 열 모두이다. 빼다가 $P$가 비면 안쪽 고리를 끝낸다(scipy의 C 판은 이때 음수 개수로 계속 돌아 첨자 배열 앞을 쓴다).
\item 반복 수는 최소제곱 풀이(안쪽 고리 한 바퀴)의 횟수이고 상한의 기본은 $3n-1$이다(scipy의 $3n$과 같은 자리에서 멈춘다). 정확한 산술의 유한 종료(\cref{thm:nnls-finite})는 부분집합의 수만큼을 허용하지만 \cref{tab:nnls-kkt}의 반복은 $n$의 $0.5$--$1.4$배였다.
\end{enumerate}

\cref{tab:nnls-kkt}는 따로 빌드한 stardust의 \code{nnls}가 돌려준 $x$에서 \cref{eq:nnls-kkt}를 잰 결과이다. 조건수를 정한 무작위 $A$에서 양수인 계수의 $|w_j|$와 $0$인 계수의 $w_j$의 양수 부분을 $\|A\|_2\|b\|$로 나누었다. 또 같은 $P$ 위의 정규방정식을 정확한 유리수로 풀어 그 해가 \cref{eq:nnls-kkt}를 정확히 만족하는지, 곧 stardust가 고른 $P$가 정답인지를 확인했다. 네 경우 모두 정답이었다. $x$의 상대오차는 $\kappa(A)$보다 $P$ 위 부분 문제의 조건수 $\kappa(A_P)$를 따랐다. 상대오차를 $\kappa(A_P)u$로 나누면 네 경우 모두 $0.25$--$0.46$이다(\code{figures/08/nnls\_condP.py}). \cref{fig:nnls-traj}는 가장 큰 문제에서 반복마다 $|P|$와 비용이 어떻게 변하는지 보인다. $|P|$는 대체로 늘지만 열을 빼는 걸음이 $10$번 있었고 비용은 단조 감소했다.

\begin{figure}[htbp]
\centering
\begin{tikzpicture}
\begin{axis}[sdplot,width=0.45\textwidth,height=0.38\textwidth,xlabel={반복},ylabel={$|P|$},xmin=0,xmax=112,ymin=0,ymax=36]
\addplot[black,const plot] table[x=k,y=P]{figures/08/nnls_traj.dat};
\end{axis}
\end{tikzpicture}\hfill
\begin{tikzpicture}
\begin{semilogyaxis}[sdplot,width=0.45\textwidth,height=0.38\textwidth,xlabel={반복},ylabel={비용$-$최종 비용},xmin=0,xmax=112,ymin=1e-4,ymax=100]
\addplot[black,mark=*,mark size=0.5pt,restrict y to domain=-20:10] table[x=k,y=cost]{figures/08/nnls_traj.dat};
\end{semilogyaxis}
\end{tikzpicture}
\caption{$(m,n)=(200,80)$, $\kappa(A)=10^8$의 비음 최소제곱을 따로 빌드한 stardust의 \code{nnls}로 풀 때 반복마다의 양수 계수 수 $|P|$(왼쪽)와 비용과 최종 비용의 차(오른쪽; \code{max\_iter}$=k$로 멈춘 반복점, \code{figures/08/nnls\_traj.py}).}
\label{fig:nnls-traj}
\end{figure}

\begin{table}[htbp]
\begin{sdtable}{colspec={cccccccc}, row{1}={font=\small\bfseries}}
\toprule
$(m,n)$ & $\kappa(A)$ & 반복 & $|P|$ & $\kappa(A_P)$ & 양수 계수의 $|w_j|$ & 활성 집합 & $x$의 상대오차 \\
\midrule
$(30,10)$ & $10$ & $5$ & $5$ & $3.7$ & $3.9\times10^{-17}$ & 정답 & $1.0\times10^{-16}$ \\
$(60,30)$ & $10^3$ & $21$ & $15$ & $59$ & $9.8\times10^{-16}$ & 정답 & $3.0\times10^{-15}$ \\
$(100,50)$ & $10^6$ & $60$ & $24$ & $2.6\times10^3$ & $3.3\times10^{-15}$ & 정답 & $1.3\times10^{-13}$ \\
$(200,80)$ & $10^8$ & $110$ & $34$ & $2.1\times10^4$ & $3.1\times10^{-14}$ & 정답 & $5.7\times10^{-13}$ \\
\bottomrule
\end{sdtable}
\caption{따로 빌드한 stardust의 \code{nnls}에서 잰 KKT 잔차(양수 계수의 $|w_j|/(\|A\|_2\|b\|)$; $0$인 계수의 $w_j$는 네 경우 모두 양수가 없었다)와 활성 집합의 정답 여부(정확한 유리수 계산; \code{figures/08/nnls\_kkt.py}).}
\label{tab:nnls-kkt}
\end{table}

상자 제약 $l\le x\le u$의 선형 최소제곱 \code{lsq\_linear}에는 같은 생각을 양쪽 경계로 넓힌 \code{'bvls'}(Stark와 Parker\cite{stark1995}: 자유 변수 위의 최소제곱을 되풀이하며 경계를 넘는 변수를 묶고, KKT 잔차가 가장 큰 변수를 하나씩 푸는 활성 집합법)와 9권에서 다룰 신뢰 영역 반사법 \code{'trf'}가 있다. \cref{tab:box-kkt}의 아래 두 줄은 같은 문제에서 두 방법의 \code{optimality}와 직접 계산한 사영 기울기 노름이다. \code{'bvls'}는 KKT 잔차가 $10^{-13}$까지 내려가 활성 집합을 정확히 찾은 반면 \code{'trf'}는 비용의 상대 변화로 멈추어 사영 기울기가 $3\times10^{-5}$에 머물렀다. \code{'trf'}의 \code{optimality} $4.5\times10^{-8}$은 경계 근처의 성분을 경계까지의 거리로 줄인 Coleman--Li 척도의 값이라 같은 점의 사영 기울기보다 훨씬 작게 나온다.

\begin{algorithm}[htbp]
\begin{algorithmic}[1]
\Require $A$, $b$, 상자 $l\le x\le u$, \code{tol}
\Ensure $x$
\State $x\gets\operatorname{lstsq}(A,b)$; \textbf{if} $l\le x\le u$ \textbf{then} \Return $x$ \Comment{경계 없는 해가 안이면 끝}
\State 밖의 성분을 경계에 묶는다(묶인 쪽의 부호 $o_j\in\{-1,+1\}$, 자유이면 $o_j=0$)
\Repeat \Comment{초기화: 자유 변수 위에서 풀고 넘는 것을 묶는다}
  \State $F\gets\{j:o_j=0\}$; $z\gets\operatorname{lstsq}(A_F,\ b-A_{F^c}x_{F^c})$; 넘는 $j\in F$를 묶고 나머지 $x_F\gets z$
\Until{넘는 것이 없다 \textbf{or} $F=\emptyset$}
\State $g\gets A^{\mathsf T}(Ax-b)$; $\mathrm{opt}\gets\max_j(o_j=0\ ?\ |g_j|:o_jg_j)$
\While{$\mathrm{opt}\ge\mathrm{tol}$ \textbf{and} 상한 전}
  \State $t\gets\arg\max_jo_jg_j$; $o_t\gets0$ \Comment{KKT를 가장 크게 어긴 묶인 변수를 푼다}
  \Loop
    \State $F\gets\{j:o_j=0\}$; $z\gets\operatorname{lstsq}(A_F,\ b-A_{F^c}x_{F^c})$
    \State \textbf{if} $l_F\le z\le u_F$ \textbf{then} $x_F\gets z$; \textbf{break}
    \State $\alpha\gets$ 경계에 처음 닿는 걸음; $x_F\gets x_F+\alpha(z-x_F)$; 닿은 변수를 묶는다
  \EndLoop
  \State 비용의 상대 감소가 \code{tol}$\cdot$비용보다 작으면 멈춘다; $g$, $\mathrm{opt}$를 갱신
\EndWhile
\State 묶인 변수는 정확히 경계 값, 나머지는 $[l,u]$로 자른다
\end{algorithmic}
\caption{경계가 있는 최소제곱 BVLS (stardust \code{linear.cpp}의 \code{bvls})}
\label{alg:bvls}
\end{algorithm}

\cref{alg:bvls}는 \cref{alg:nnls}와 같은 뼈대를 양쪽 경계로 넓힌 것이다. 9행은 묶인 쪽에서 안으로 들어갈 때 목적 함수가 줄어드는 변수 가운데 가장 크게 줄어드는 변수를 고르고(\cref{eq:nnls-kkt}의 상자판) 10--14행은 Lawson--Hanson의 안쪽 고리 그대로이다. 비용은 다르다. \code{nnls}는 Householder QR을 열 하나씩 고쳐 가며 바깥 고리 한 번에 $O(mn)$이지만 \code{bvls}는 scipy처럼 자유 열에 대한 \code{lstsq}(LAPACK \code{gelsd}, $\mathrm{rcond}=\epsilon\max(m,|F|)$)를 처음부터 다시 하므로 한 번에 $O(m|F|^2)$이다. 마지막 줄에서 stardust는 scipy와 달리 반올림 때문에 경계를 아주 조금 넘은 자유 변수를 경계로 잘라 결과가 언제나 상자 안에 있게 한다.

\section{Smoothness and Gradient Methods}
\label{sec:smooth}

\subsection{Gradient Descent}

가장 단순한 알고리즘은 기울기의 반대 방향으로 가는 \term{경사 하강법}{gradient descent} $x_{k+1}=x_k-\alpha\nabla f(x_k)$이다. Cauchy가 1847년에 연립방정식을 풀려고 제안했다\cite{cauchy1847,lemarechal2012cauchy}. \cref{prop:descent-lemma}에 $y=x-\alpha\nabla f(x)$를 넣으면 $f(y)\le f(x)-\alpha(1-L\alpha/2)\|\nabla f(x)\|^2$이므로 $\alpha=1/L$에서 걸음마다
\begin{equation}
\label{eq:sufficient-decrease}
f(x_{k+1})\le f(x_k)-\frac1{2L}\|\nabla f(x_k)\|^2
\end{equation}
만큼 함수값이 준다.

\begin{theorem}
\label{thm:gd-rates}
$f$가 $L$-매끄러운 볼록 함수이고 최소점 $x_\ast$가 있으면 $\alpha=1/L$의 경사 하강법은
\[
f(x_k)-f(x_\ast)\le\frac{L\|x_0-x_\ast\|^2}{2k}
\]
를 만족한다. $f$가 $\mu$-강볼록이기도 하면 $f(x_k)-f(x_\ast)\le(1-\mu/L)^k\bigl(f(x_0)-f(x_\ast)\bigr)$이다.
\end{theorem}

\begin{proof}
볼록성에서 $f(x_k)-f(x_\ast)\le\nabla f(x_k)^{\mathsf T}(x_k-x_\ast)$이다. 이것과 \cref{eq:sufficient-decrease}를 합치고 $g=\nabla f(x_k)$로 쓰면
\[
f(x_{k+1})-f(x_\ast)\le g^{\mathsf T}(x_k-x_\ast)-\frac1{2L}\|g\|^2=\frac L2\Bigl(\|x_k-x_\ast\|^2-\|x_{k+1}-x_\ast\|^2\Bigr)
\]
이다(가운데 식을 완전제곱으로 정리한다). $k=0,\dots,K-1$에 대해 더하면 오른쪽이 망원합이 되어 $\sum_{k=1}^K(f(x_k)-f(x_\ast))\le(L/2)\|x_0-x_\ast\|^2$이고 \cref{eq:sufficient-decrease}에서 $f(x_k)$가 감소하므로 좌변은 $K(f(x_K)-f(x_\ast))$ 이상이다. 강볼록이면 \cref{prop:descent-lemma}의 $\|\nabla f(x)\|^2\ge2\mu(f(x)-f(x_\ast))$를 \cref{eq:sufficient-decrease}에 넣어 $f(x_{k+1})-f(x_\ast)\le(1-\mu/L)(f(x_k)-f(x_\ast))$를 얻는다.
\end{proof}

강볼록인 경우 오차를 $\epsilon$배로 줄이려면 대략 $\kappa\ln(1/\epsilon)$걸음이 든다. 걸음을 $1/L$로 고정하지 않고 방향을 따라 정확히 최소화해도 사정은 크게 나아지지 않는다. 이차 함수에서 정확한 선 탐색의 경사 하강은 함수값의 오차를 걸음마다 $((\kappa-1)/(\kappa+1))^2$배 이하로 줄이고 이 비는 최악의 출발점에서 실제로 이루어진다. 이 비율도 볼록성으로 증명한다.

\begin{proposition}
\label{prop:kantorovich}
$A$가 고유값이 $[m,M]$($0<m\le M$)에 든 대칭 양정치 행렬이면 모든 $d\ne0$에서
\[
(d^{\mathsf T}Ad)(d^{\mathsf T}A^{-1}d)\le\frac{(m+M)^2}{4mM}\|d\|^4
\]
이다(Kantorovich 부등식). 따라서 $f=(1/2)x^{\mathsf T}Ax-b^{\mathsf T}x$의 정확한 선 탐색 경사 하강은 $f(x_{k+1})-f_\ast\le((M-m)/(M+m))^2(f(x_k)-f_\ast)$를 만족한다.
\end{proposition}

\begin{proof}
$\|d\|=1$로 두고 $A$의 고유 기저에서 $p_i=d_i^2$($\sum p_i=1$)로 쓰면 $t=d^{\mathsf T}Ad=\sum p_i\lambda_i$, $d^{\mathsf T}A^{-1}d=\sum p_i/\lambda_i$이다. $1/\lambda$는 $[m,M]$에서 볼록이므로 현 아래에 있어 $1/\lambda\le(m+M-\lambda)/(mM)$이고 $p_i$로 평균하면 $d^{\mathsf T}A^{-1}d\le(m+M-t)/(mM)$이다. 따라서 곱은 $t(m+M-t)/(mM)\le(m+M)^2/(4mM)$이다(이차식의 최댓값). 정확한 선 탐색의 걸음은 $\alpha=g^{\mathsf T}g/g^{\mathsf T}Ag$이고 $f(x)-f_\ast=(1/2)g^{\mathsf T}A^{-1}g$이므로 계산하면
\[
f(x_{k+1})-f_\ast=\Bigl(1-\frac{(g^{\mathsf T}g)^2}{(g^{\mathsf T}Ag)(g^{\mathsf T}A^{-1}g)}\Bigr)(f(x_k)-f_\ast)
\]
이고 괄호 안은 부등식에서 $1-4mM/(m+M)^2=((M-m)/(M+m))^2$ 이하이다.
\end{proof}

\cref{fig:zigzag}가 그 모습이다. 방향이 등고선에 수직이라 긴 골짜기를 가로질러 이리저리 튀며 느리게 내려간다.

\begin{figure}[htbp]
\centering
\begin{tikzpicture}
\begin{axis}[sdplot,xlabel={$x_1$},ylabel={$x_2$},xmin=-1,xmax=11,ymin=-1.6,ymax=1.6,width=0.8\textwidth,height=0.36\textwidth]
\foreach \c in {1,5,15,30,55} {\addplot[black!35,domain=0:360,samples=90,variable=\t] ({sqrt(2*\c)*cos(\t)},{sqrt(2*\c/10)*sin(\t)});}
\addplot[black,thick,mark=*,mark size=1pt] table[x=x,y=y]{figures/08/gd_exact.dat};
\addlegendentry{정확한 선 탐색}
\addplot[black,dashed,mark=o,mark size=1.2pt] table[x=x,y=y]{figures/08/gd_fixed.dat};
\addlegendentry{걸음 $1/L$}
\end{axis}
\end{tikzpicture}
\caption{$f=(x_1^2+10x_2^2)/2$($\kappa=10$)에서 $(10,1)$을 출발한 경사 하강의 처음 열두 걸음(회색은 등고선; \code{figures/08/zigzag.py}). 정확한 선 탐색은 함수값을 걸음마다 정확히 $(9/11)^2\approx0.669$배로 줄이며 지그재그로 가고, 걸음 $1/L$은 첫걸음에 $x_2$를 없앤 뒤 $x_1$ 방향으로 $0.9$배씩 느리게 간다.}
\label{fig:zigzag}
\end{figure}

곡률 정보를 쓰면 이 의존성이 줄어든다. 이차 함수에서 켤레 기울기법은 $\sqrt\kappa\ln(1/\epsilon)$에 비례하는 걸음이면 충분하고(63장) 준뉴턴법은 곡률을 쌓아 가며 결국 초선형으로 수렴한다(9권). \cref{fig:kappa}는 고유값이 $[1,\kappa]$에 고르게 퍼진 $100$차원 이차 함수에서 함수값의 오차가 $10^{-10}$배가 될 때까지의 걸음 수이다. 경사 하강과 켤레 기울기법은 여기서 짠 것이고 BFGS와 L-BFGS-B는 따로 빌드한 stardust의 \code{minimize}(\code{max\_iter}를 바꿔 가며 처음 닿는 반복을 찾았다)이다. $\kappa=10^4$에서 경사 하강은 $40412$걸음, 켤레 기울기법은 $54$걸음, stardust의 BFGS는 $108$걸음, L-BFGS-B는 $70$걸음이었다. 경사 하강의 걸음 수는 \cref{thm:gd-rates}의 상한 $\kappa\ln10^{10}$보다 훨씬 적다. $\kappa=10^4$에서 약 $0.18$배이고 이 비는 $\kappa$가 커질수록 줄어든다($\kappa=3$에서 $0.40$). 이차 함수에서는 고유값마다 함수값의 오차가 $(1-\lambda_i/L)^{2k}$로 줄어 지수가 상한의 두 배이다. 또 가장 느린 방향이 처음 함수값에서 차지하는 몫이 $1/(n\kappa)$ 정도로 작아서 그 방향은 $10^{-10}$배까지 줄일 필요가 없다. 켤레 기울기법이 $54$걸음에서 멈춘 것은 유한 종료 때문이다. $n=100$차원에서는 많아야 $n$걸음이면 끝나며 $\sqrt\kappa$ 법칙과는 관계가 없다($\sqrt\kappa\ln10^{10}$의 어림으로는 $\kappa=10^4$에서 약 $600$걸음이다). BFGS는 $n$보다 많은 $108$걸음을 썼다. Wolfe 선 탐색이 정확하지 않아서 이차 함수에서도 유한 종료가 성립하지 않는다.

\begin{figure}[htbp]
\centering
\begin{tikzpicture}
\begin{loglogaxis}[sdplot,xlabel={조건수 $\kappa$},ylabel={걸음 수},xmin=2,xmax=1.5e4,ymin=5,ymax=3e5]
\addplot[black,thick,mark=*,mark size=1.1pt] table[x=kappa,y=gd]{figures/08/kappa.dat};
\addlegendentry{경사 하강 ($1/L$)}
\addplot[black!50,dotted,thick] table[x=kappa,y=bound]{figures/08/kappa.dat};
\addlegendentry{\cref{thm:gd-rates}의 상한}
\addplot[black,dashed,mark=square*,mark size=1.1pt] table[x=kappa,y=bfgs]{figures/08/kappa.dat};
\addlegendentry{stardust BFGS}
\addplot[black,mark=triangle*,mark size=1.4pt] table[x=kappa,y=lbfgsb]{figures/08/kappa.dat};
\addlegendentry{stardust L-BFGS-B}
\addplot[black!60,mark=o,mark size=1.3pt] table[x=kappa,y=cg]{figures/08/kappa.dat};
\addlegendentry{켤레 기울기법}
\end{loglogaxis}
\end{tikzpicture}
\caption{함수값의 오차가 처음의 $10^{-10}$배가 될 때까지의 걸음 수($n=100$의 대각 이차 함수, 고유값은 $[1,\kappa]$에 고르게; \code{figures/08/kappa.py}). 측정에서 BFGS는 \code{gtol}$=0$, L-BFGS-B는 \code{gtol}$=0$과 \code{ftol}$=0$으로 두어 걸음 수만 보았다. 경사 하강은 $\kappa$에 비례하고 켤레 기울기법은 $n=100$의 유한 종료에 막혀 그 아래에 머문다.}
\label{fig:kappa}
\end{figure}

\subsection{Projected Gradient}

상자 같은 단순한 집합 $C$ 위에서는 걸음마다 사영하는 \term{사영 경사 하강법}{projected gradient method} $x_{k+1}=P_C(x_k-\alpha\nabla f(x_k))$이 자연스럽다. \cref{sec:kkt}에서 본 대로 그 고정점은 KKT 점과 같다.

\begin{proposition}
\label{prop:projected-gd}
$C$가 공집합이 아닌 닫힌 볼록 집합이고 $f$가 두 번 연속 미분 가능하며 $\mu I\preceq\nabla^2f\preceq LI$이면 $\alpha=1/L$의 사영 경사 하강법은 $C$ 위의 최소점 $x_\ast$에 대해 $\|x_{k+1}-x_\ast\|\le(1-\mu/L)\|x_k-x_\ast\|$를 만족한다.
\end{proposition}

\begin{proof}
$x_\ast=P_C(x_\ast-\alpha\nabla f(x_\ast))$이고 \cref{thm:projection}에서 $P_C$는 비확대이므로 $\|x_{k+1}-x_\ast\|\le\|x_k-x_\ast-\alpha(\nabla f(x_k)-\nabla f(x_\ast))\|$이다. 적분 꼴의 평균값 정리에서 $\nabla f(x_k)-\nabla f(x_\ast)=\bar H(x_k-x_\ast)$, $\bar H=\int_0^1\nabla^2f(x_\ast+t(x_k-x_\ast))\,dt$이고 $\mu I\preceq\bar H\preceq LI$이므로 $\|I-\alpha\bar H\|_2\le\max(|1-\alpha\mu|,|1-\alpha L|)=1-\mu/L$이다.
\end{proof}

$L$을 모르면 걸음을 되짚어 줄이는 판(\cref{alg:pgd})을 쓴다. 하강 보조정리에서 $\alpha\le1/L$이면 충분 감소 조건이 성립하므로 되짚기는 $\alpha\ge\beta/L$에서 반드시 끝난다. 반복 하나의 비용은 사영의 $O(n)$에 함수와 기울기의 계산이 전부이다.

\begin{algorithm}[htbp]
\begin{algorithmic}[1]
\Require $x_0\in C$, 처음 걸음 $\alpha_0$, $\beta,c\in(0,1)$, 문턱 \code{pgtol}
\Ensure KKT 점의 근사 $x$
\For{$k=0,1,\dots$}
  \State $g\gets\nabla f(x_k)$; \textbf{if} $\|x_k-P_C(x_k-g)\|_\infty\le\mathrm{pgtol}$ \textbf{then} \Return $x_k$ \Comment{\cref{alg:projgr}}
  \State $\alpha\gets\alpha_0$
  \Repeat
    \State $x^+\gets P_C(x_k-\alpha g)$; $\alpha\gets\beta\alpha$
  \Until{$f(x^+)\le f(x_k)+c\,g^{\mathsf T}(x^+-x_k)$} \Comment{사영 경로 위의 충분 감소}
  \State $x_{k+1}\gets x^+$
\EndFor
\end{algorithmic}
\caption{되짚기 걸음의 사영 경사 하강}
\label{alg:pgd}
\end{algorithm}

사영 경사 하강은 상자 경계를 다루는 방법으로는 옳지만 느리다. L-BFGS-B\cite{byrd1995lbfgsb}는 사영 경로를 활성 집합을 고르는 데만 쓴다. 기울기 방향의 사영 경로 $x(t)=P_C(x-tg)$는 $t$가 커지면서 성분이 하나씩 경계에 닿아 멈추는 꺾인 선이다. 성분 $i$가 경계에 닿는 시각은 $g_i<0$이면 $t_i=(x_i-u_i)/g_i$, $g_i>0$이면 $(x_i-l_i)/g_i$이고 그사이의 구간에서 경로는 직선이다. 이 경로를 따라 이차 모형 $m(y)=f(x)+g^{\mathsf T}(y-x)+(1/2)(y-x)^{\mathsf T}B(y-x)$($B$는 준뉴턴 행렬)의 첫 국소 최소점을 \term{일반화 Cauchy 점}{generalized Cauchy point}이라 한다(\cref{alg:cauchy}). 그 점에서 경계에 닿은 성분을 활성 집합으로 고정하고 나머지 자유 변수 위에서만 준뉴턴 걸음을 푼다(9권).

\begin{algorithm}[htbp]
\begin{algorithmic}[1]
\Require $x\in C$, $g$, 모형의 곡률 $B$, 상자 $l,u$
\Ensure 일반화 Cauchy 점 $x^c$
\State 꺾이는 시각 $t_i$를 계산한다; $d\gets-g$ (경계에 붙어 밖을 향하는 성분, $l_i=u_i$, $g_i=0$인 성분은 $d_i\gets0$)
\State $t_{\mathrm{old}}\gets0$; $f'\gets g^{\mathsf T}d$; $f''\gets d^{\mathsf T}Bd$; $\Delta t_{\min}\gets-f'/f''$ \Comment{첫 구간의 직선에서 모형의 최소점}
\While{남은 꺾임이 있다}
  \State $t_{\mathrm{next}}\gets$ 다음 꺾임 \Comment{첫 번째는 훑어서, 그다음부터 힙에서}
  \State \textbf{if} $\Delta t_{\min}<t_{\mathrm{next}}-t_{\mathrm{old}}$ \textbf{then} \textbf{break} \Comment{구간 안에서 멈춘다}
  \State 성분 $b$를 경계에 고정하고 $d_b\gets0$; $f'$와 $f''$를 그 성분만큼 고친다; $f''\gets\max(\epsilon f''_0,f'')$
  \State $t_{\mathrm{old}}\gets t_{\mathrm{next}}$; $\Delta t_{\min}\gets-f'/f''$ ($d=0$이면 $0$) \Comment{다음 구간}
\EndWhile
\State $x^c\gets x(t_{\mathrm{old}}+\max(0,\Delta t_{\min}))$
\end{algorithmic}
\caption{사영 경로 위의 일반화 Cauchy 점 (stardust \code{lbfgsb.cpp}의 \code{cauchy}, 힙은 \code{hpsolb})}
\label{alg:cauchy}
\end{algorithm}

이 알고리즘은 구간을 넘어갈 때 $f'$와 $f''$를 처음부터 다시 계산하지 않고 경계에 닿은 성분 하나의 기여만 빼고 더한다. $B$가 L-BFGS의 압축된 꼴(단위 행렬의 배수에 계수 $2m$의 보정)이므로 그 갱신은 $O(m^2)$이고(중간 행렬 $2m\times2m$의 곱과 삼각 풀이, \code{bmv}) 구간 하나의 비용이 $n$에 의존하지 않는다. 시작할 때 $p=W^{\mathsf T}d$를 만드는 $O(mn)$이 따로 든다. 꺾이는 시각을 처음에 모두 정렬하면 $O(n\log n)$이지만 대개는 몇 구간만에 멈추므로 첫 꺾임은 최솟값 하나를 훑어 찾고 그 꺾임을 지난 뒤에야 힙을 만들어(\code{hpsolb}, 삽입식이라 최악 $O(n\log n)$) 하나씩 $O(\log n)$에 꺼낸다. 경계 경우는 세 가지를 따로 다룬다. 곡률 $f''$에는 처음 값의 $\epsilon$배라는 바닥을 두어 그 아래로 내려가지 않게 한다. 모든 성분이 경계에 닿으면 마지막 걸음이 $0$이고 경계 없는 성분이 남으면 마지막 반직선 위의 최소점으로 간다. ``마지막 꺾임'' 판정은 scipy의 C 판에서 첨자를 $0$부터 세게 바꾼 뒤로 결코 참이 되지 않는데 stardust도 이 판정을 그대로 두었다.

\section{Limits of Finite Precision}
\label{sec:opt-precision}

최솟값 근처에서 함수는 평평하다. $f(x_\ast+h)\approx f(x_\ast)+(1/2)f''(x_\ast)h^2$이므로 $(1/2)|f''|h^2$이 $u|f(x_\ast)|$보다 작으면 binary64로 계산한 함수값은 $f(x_\ast)$와 구별되지 않는다. 함수값만 보고 최소점의 위치를 찾으면 그 정확도는 대략
\[
|h|\lesssim\sqrt{\frac{2u|f(x_\ast)|}{|f''(x_\ast)|}}
\]
에 머물고 상대 정확도로는 $\sqrt u\approx10^{-8}$ 수준이다. $f(x)=(x-1)^2+1$을 binary64로 계산하면 $|h|\le1.05\times10^{-8}$인 모든 $x=1+h$에서 $f(x)$가 정확히 $1$이다(\cref{fig:flat}; $1$ 위쪽으로 이웃한 double $4.7\times10^7$개, 아래쪽으로 $9.5\times10^7$개). 74장의 앞쪽 차분의 정확도가 $\sqrt u$에 머문 것과 모양이 닮았다. 거기서는 절단 오차와 반올림 오차의 균형이, 여기서는 최소점 근처의 이차적인 평평함이 제곱근을 낳는다.

\begin{figure}[htbp]
\centering
\begin{tikzpicture}
\begin{axis}[sdplot,xlabel={$h/10^{-8}$},ylabel={$(\mathrm{fl}(f(1+h))-1)/2^{-52}$},xmin=-5,xmax=5,ymin=-0.3,ymax=12]
\addplot[black,const plot,thick] table[x expr=\thisrow{h}*1e8,y expr=\thisrow{fm1}*4503599627370496]{figures/08/flat.dat};
\addlegendentry{binary64}
\addplot[black!50,dashed,domain=-5:5,samples=100] {x^2*0.45035996};
\addlegendentry{$h^2$}
\end{axis}
\end{tikzpicture}
\caption{$f(x)=(x-1)^2+1$을 binary64로 계산한 값. 최소점 둘레의 너비 약 $2.1\times10^{-8}$에서 값이 정확히 $1$이고 그 밖에서는 $2u=2^{-52}$ 간격의 계단이다(\code{figures/08/precision.py}).}
\label{fig:flat}
\end{figure}

기울기를 쓰면 사정이 낫다. $\nabla f(x_\ast+h)\approx\nabla^2f\,h$는 $h$에 비례하므로 기울기가 상대오차 $u$로 계산되는 한 위치를 $u$에 가까운 정확도로 정할 수 있다. 따로 빌드한 stardust로 $f(x)=e^x-2x$(최소점 $\ln2$)를 풀어 보면 차이가 뚜렷하다. 함수값만 쓰는 \code{minimize\_scalar}의 \code{'brent'}는 상대오차 $1.2\times10^{-9}$, \code{'golden'}은 $2.1\times10^{-9}$이고 기울기를 쓰는 \code{minimize}의 BFGS(\code{gtol}$=10^{-15}$)는 $1.6\times10^{-16}$이었다. 이 함수의 평평한 구간은 $\ln2$ 위로 $6.8\times10^{-9}$이므로 앞의 둘은 한계에 닿은 것이다. \code{'bounded'}는 scipy의 fminbound를 따라 허용오차를 $\sqrt{\epsilon}\,|x|+\mathrm{xatol}/3$($\epsilon=2u$는 기계 엡실론)으로 잡는데 \code{xatol}의 기본값 $10^{-5}$가 이를 지배해 상대오차가 $5.4\times10^{-7}$이었다. \code{xatol}을 $10^{-15}$로 내리면 $1.4\times10^{-9}$로 같은 한계에 닿았다. $(x-1)^2+1$에서는 세 방법이 모두 $x=1$을 정확히 돌려주었지만 그 이유는 평평함과 상관이 없다. \code{'brent'}와 \code{'golden'}은 기본 구간 탐색이 $0$과 $1$에서 시작해 처음부터 $x=1$을 짚고 그보다 엄격히 작은 값을 끝내 만나지 못하기 때문이고 \code{'bounded'}는 포물선 보간의 꼭짓점이 반올림 없이 계산되었기 때문이다. 꼭짓점을 $c$로 옮긴 $(x-c)^2+1$에서는 결과가 경우마다 다르다. $c=0.7$에서 \code{'brent'}는 $-1.5\times10^{-8}$로 평평한 구간의 한계에 머물렀고 $c=1.3$, $2.9$에서는 반올림 수준으로 맞았다(\code{figures/08/quad\_exact.py}).

함수값의 감소로 멈추는 판정도 같은 한계를 만난다. \cref{tab:box-kkt}의 위 네 줄에서는 $50$차원의 상자 제약 이차 문제를 stardust의 L-BFGS-B로 풀었다. 기본 옵션에서는 상대 감소 판정이 사영 기울기 $1.0\times10^{-3}$에서 먼저 걸린다(\cref{sec:optimality}). 둘째 줄부터는 \code{ftol}$=0$으로 그 판정을 끄고 기울기 판정만 보았다. \code{gtol}을 $10^{-5}$에서 $10^{-11}$까지 줄여도 사영 기울기는 $8.6\times10^{-7}$에서 더 내려가지 않고 멈춤 사유는 ``함수값의 상대 감소가 \code{factr}$\cdot\epsilon$ 이하''였다(\code{ftol}$=0$으로 둔 실험이므로 이 사유가 나왔다면 함수값이 한 걸음에서 전혀 줄지 않았다). 반복이 함수값을 더 줄이지 못하는 \cref{fig:flat}의 평평한 구간에 닿았다. 이때의 기울기는 $\sqrt{2u|f|\lambda}$ 정도($\lambda$는 남은 방향의 곡률)로 어림된다. 이 문제에서 $|f|=328$이고 자유 변수 부분의 고유값이 $[7.1,96.6]$이므로 어림은 $0.7$--$2.7\times10^{-6}$이고 잰 값 $8.6\times10^{-7}$이 그 범위에 든다.

\begin{table}[htbp]
\begin{sdtable}{colspec={X[1.6]X[1.2]cccc}, row{1}={font=\small\bfseries}}
\toprule
방법 & 판정 & 반복 & 활성 & 사영 기울기 & 보고한 \code{optimality} \\
\midrule
L-BFGS-B, 기본 \code{ftol} & 상대 감소 & $19$ & $10$ & $1.0\times10^{-3}$ & -- \\
L-BFGS-B, \code{gtol}$=10^{-5}$ & 사영 기울기 & $28$ & $10$ & $8.1\times10^{-6}$ & -- \\
L-BFGS-B, \code{gtol}$=10^{-8}$ & 상대 감소 & $33$ & $10$ & $8.6\times10^{-7}$ & -- \\
L-BFGS-B, \code{gtol}$=10^{-11}$ & 상대 감소 & $33$ & $10$ & $8.6\times10^{-7}$ & -- \\
\code{lsq\_linear}, \code{'trf'} & 비용의 변화 & $12$ & $29$ & $3.3\times10^{-5}$ & $4.5\times10^{-8}$ \\
\code{lsq\_linear}, \code{'bvls'} & KKT 잔차 & $9$ & $29$ & $1.4\times10^{-13}$ & $1.4\times10^{-13}$ \\
\bottomrule
\end{sdtable}
\caption{상자 제약 문제에서 따로 빌드한 stardust가 멈춘 점의 사영 기울기 노름(\cref{def:projected-gradient}, 직접 계산). 위 네 줄은 $n=50$, 고유값 $[1,100]$의 이차 함수와 $[-1,1]^{50}$이고 첫 줄은 기본 옵션(\code{gtol}을 $10^{-5}$에서 $10^{-11}$까지 바꿔도 같음), 그 아래 셋은 \code{ftol}$=0$이다. 마지막 두 줄은 $80\times50$의 최소제곱과 $[0,1]^{50}$이다(\code{figures/08/box\_kkt.py}, \code{box\_ftol.py}).}
\label{tab:box-kkt}
\end{table}

\section{Testing Optimality Code}
\label{sec:opt-testing}

최적화 코드는 참값을 모르므로 시험하기 어렵다. 그 빈자리를 이 장의 이론이 메운다. 첫째, KKT 조건이나 쌍대 간격처럼 답을 증명하는 양을 독립적으로 계산한다. \cref{tab:nnls-kkt}에서는 활성 집합을 정확한 유리수로 검증했고 \cref{fig:lasso-gap}에서는 간격을 유리수로 계산했다. 이렇게 하면 계산한 답이 맞는지와 함께 판정 자체의 계산(간격의 소거 같은)을 믿을 만한지도 드러난다. 둘째, 닫힌 꼴의 해가 있는 문제를 쓴다. 이차 함수, 하나뿐인 열의 라소(\cref{prop:soft}의 연성 임계), 상자 안의 분리 가능한 문제가 그렇다. 셋째, 같은 문제를 다른 방법으로 풀어 견준다(\code{'bvls'}와 \code{'trf'}, NNLS와 \code{lsq\_linear}의 $[0,\infty)$). 넷째, 멈춤 사유를 확인한다. \cref{tab:box-kkt}처럼 \code{converged}가 참이어도 멈춘 이유가 기울기가 아니라 함수값의 정체일 수 있다.

\cref{tab:opt-tests}에 이 장에서 쓴 검증과 그 결과를 모았다.

\begin{table}[htbp]
\begin{sdtable}{colspec={X[1.6]X[2.6]X[2]}, row{1}={font=\small\bfseries}}
\toprule
대상 & 무엇과 견주었나 & 결과 \\
\midrule
\code{nnls} & 활성 집합 위의 정규방정식을 유리수로 풀어 KKT 확인 & 네 경우 모두 정답, 잔차 $3\times10^{-14}$ 이하 \\
\code{lasso}의 간격 & 같은 점의 간격을 유리수로 계산 & 보고한 $3.75\times10^{-12}$ 대 정확한 $3.72\times10^{-12}$ \\
간격의 식 & 정의대로의 $P-D$ 대 \cref{eq:gap-stable} & 바닥 $4\times10^{-13}$ 대 $4\times10^{-15}$ \\
L-BFGS-B, \code{lsq\_linear} & 직접 계산한 사영 기울기 & 멈춤 사유에 따라 $10^{-5}$대에서 $10^{-13}$ \\
\code{minimize\_scalar} & 평평한 구간의 너비 & 평평한 구간 안에서 멈춤(상대오차 $10^{-9}$대) \\
\bottomrule
\end{sdtable}
\caption{이 장의 검증 요약(따로 빌드한 stardust; 스크립트는 \code{figures/08/}).}
\label{tab:opt-tests}
\end{table}

stardust 자신의 시험은 scipy 1.18.1을 기준으로 삼는다. 설계 문서에 따르면 BFGS를 뺀 \code{minimize}의 다섯 방법은 성공 경로에서 사용자 함수를 부르는 점이 scipy와 비트까지 같고(약속은 아니고 지금의 사실이다) BFGS는 역 Hessian 갱신의 계산 순서를 바꾸어 반올림만큼 다르다. 시험은 이 호출점들을 기록해 견준다. \code{nnls}와 \code{lsq\_linear}은 scipy가 부르는 BLAS·LAPACK 호출과 작업 배열의 배치까지 같게 두어 결과를 비트로 맞춘다. 라소는 scikit-learn과 같은 목적 함수와 간격 판정을 쓰므로 그 결과와 견준다. 또 KKT 다듬기의 결과가 경로와 한 번 풀기에서 같은 비트인지 본다.

\section{APIs in This Chapter}
\label{sec:convexity-apis}

\cref{tab:convexity-apis}에 이 장의 이론이 구현으로 나타나는 stardust API와 그것을 설명한 곳을 모았다. 각 API의 인자와 알고리즘 전체는 95장(한 변수 최소화), 97장(L-BFGS-B와 \code{minimize}), 101장(\code{lsq\_linear}, \code{nnls}), 89장(\code{logsumexp}), 126장(벌점 회귀)에서 다룬다.

\begin{table}[htbp]
\begin{sdtable}{colspec={X[2]X[3]X[1.5]}, row{1}={font=\small\bfseries}}
\toprule
API & 이 장에서 다룬 것 & 설명한 곳 \\
\midrule
\code{optimize.minimize} (\code{converged}) & 정류점 판정과 그 척도 & \cref{sec:optimality} \\
\code{optimize.minimize}, \code{'l-bfgs-b'} & 사영 기울기 노름, 일반화 Cauchy 점, 두 멈춤 판정 & \cref{alg:projgr,alg:cauchy}, \cref{tab:stop-rules,tab:box-kkt} \\
\code{optimize.nnls} & KKT에서 유도한 활성 집합법, 들이기 판정 & \cref{alg:nnls} \\
\code{optimize.lsq\_linear} & 상자 KKT, \code{optimality}의 척도, BVLS & \cref{alg:bvls}, \cref{sec:kkt} \\
\code{optimize.}\allowbreak\code{minimize\_scalar} & 함수값만으로 얻는 $\sqrt u$ 한계, \code{'bounded'}의 허용오차 & \cref{sec:opt-precision} \\
\code{regression.lasso}, \code{elastic\_net} & 부분기울기 조건, 연성 임계, 쌍대 간격, KKT 다듬기 & \cref{alg:cd,alg:cd-outer,alg:gap,alg:polish} \\
\code{special.logsumexp} & 볼록성, 최댓값 이동과 TwoSum 보정 & \cref{alg:lse} \\
\bottomrule
\end{sdtable}
\caption{이 장의 이론을 쓰는 stardust API.}
\label{tab:convexity-apis}
\end{table}

\section{Concluding Remarks}

최적화의 이론에는 두 줄기가 있다. 하나는 최소점에서 기울기와 Hessian, 제약이 있으면 승수가 만족하는 조건이다. 다른 하나는 알고리즘이 그 점에 다가가는 속도이고 이 속도는 매끄러움과 강볼록성이 정한다. 볼록성은 두 줄기를 모두 단순하게 만든다. 국소 최소와 전역 최소의 구별이 사라지고 KKT 조건이 충분조건이 되며 쌍대 간격이 최적값까지의 거리를 증명한다. 구현에서는 이 조건들이 그대로 멈춤 판정이 되었다. 사영 기울기는 KKT 잔차이다. Lawson--Hanson의 활성 집합법은 KKT 조건을 하나씩 만족시켜 가는 절차이고 라소의 간격은 쌍대 실현 가능 점 하나로 계산하는 증명서이다. 판정을 계산하는 방식도 정확도를 좌우했다. 간격은 소거 없는 꼴로 써야 바닥 아래를 읽을 수 있고 함수값만으로는 위치를 $\sqrt u$보다 정확히 정할 수 없다. 볼록 해석의 표준 참고서는 Rockafellar\cite{rockafellar1970}이고 볼록 최적화와 쌍대는 Boyd와 Vandenberghe\cite{boyd2004convex}, 일차 방법의 이론은 Nesterov\cite{nesterov2018}와 Beck\cite{beck2017}, 알고리즘 전반은 Nocedal과 Wright\cite{nocedal2006}이 자세하다. 다음 장에서는 확률론을 측도론 위에 다시 세운다.

\chapterreferences
